\begin{center}
{\Large \bfseries \textcolor{lyred}{第九章 多元函数微分法及其应用}}
\end{center}

\vspace{6pt}
{\large \bfseries \textcolor{lyred}{第9.1 节 多元函数的基本概念}}

\vspace{4pt}
{\bfseries \textcolor{lyred}{一.平面点集 n 维空间}}

{\bfseries \textcolor{lyred}{1.平面点集}}

建立了直角坐标系之后,平面可以表示为
(
)
\{
\}
,
: ,
x y
x y
=
\in 
\mathbb{R}
\mathbb{R},称为\textcolor{lyred}{坐标平面}. 
设
(
)
,
P x
y
\in \mathbb{R},记
(
)\{
\}
(
)
(
)
(
)
\{
\}
,
:
,
:
U P
P PP
x y
x
x
y
y
\delta 
\delta 
\delta 
=
<
=
-
+
-
<
, 
称为
0P 的\delta \textcolor{lyred}{邻域},简记为
(
)
U P ; 
记
(
)\{
\}
,
:0
U P
P
PP
\delta 
\delta 
=
<
<
\circ 
,称为
0P 的\textcolor{lyred}{去心}\delta \textcolor{lyred}{邻域},简记为
(
)
U P
\circ 
. 
\textcolor{lyred}{定义}.设
E \subset \mathbb{R},
0P \in \mathbb{R},(1)若
(
)
U P
\exists 
,使
(
)
U P
E
\subset 
,则称
0P 为E 的\textcolor{lyred}{内点}; 
(2)若
(
)
U P
\exists 
,使
(
)
U P
E
\cap 
=\emptyset ,则称
0P 为E 的\textcolor{lyred}{外点}; 
(3)若
0P 既不是E 的内点,也不是外点,则称
0P 为E 的\textcolor{lyred}{边界点},其全体构成的点集 
称为E 的\textcolor{lyred}{边界},记为E
\partial ; 
(4)若
(
)
U P
\forall 
\circ 
,均有
(
)
U P
E
\cap 
\neq \emptyset 
\circ 
,则称
0P 是E 的\textcolor{lyred}{聚点}. 
\textcolor{lyred}{定义}.设
E \subset \mathbb{R},(1)若E
E
\cap \partial 
=\emptyset ,则称E 为\textcolor{lyred}{开集}; 
(2)若E
E
\partial 
,则称E 为\textcolor{lyred}{闭集}. 
\textcolor{lyred}{例}.
(
)
\{
\}
,
:1
E
x y
x
y
=
<
+
<
是开集,
(
)
\{
\}
,
:1
E
x y
x
y
=
\le 
+
\le 
是闭集, 
(
)
\{
\}
,
:1
E
x y
x
y
=
<
+
\le 
即非开集,也非闭集. 
\textcolor{lyred}{定义}.设
E \subset \mathbb{R},若E 中任何两点均可以用完全包含于E 中的折线连结起来,则 
称E 是\textcolor{lyred}{连通集}. 
\textcolor{lyred}{定义}.连通的开集称为\textcolor{lyred}{(开)区域};连通的闭集称为\textcolor{lyred}{闭区域}. 
\textcolor{lyred}{定义}.设
E \subset \mathbb{R},若存在
R >
,使得
(
)
\{
\}
,
:
E
x y
x
y
R
\subset 
+
<
,则称E 为\textcolor{lyred}{有界集}; 
否则称为\textcolor{lyred}{无界集}. 
{\bfseries \textcolor{lyred}{2.n 维空间}}

\textcolor{lyred}{定义}.记
(
)
\{
\}
,
,
,
:
n
n
i
x x
x
x
=
\in 
\Lambda 
\mathbb{R}
\mathbb{R},称为\textcolor{lyred}{n 维空间},其中的元素
(
)
,
,
,
n
x
x x
x
=
\Lambda 
称为n 维空间中的\textcolor{lyred}{点},或者\textcolor{lyred}{n 维向量}. 
对于n 维向量,也可以定义加法,数乘,数量积等运算. 
\textcolor{lyred}{定义}.设
(
)
,
,
,
n
x
x x
x
=
\Lambda 
,记
n
x
x
x
x
=
+
+
+
\Lambda 
,称为x 的\textcolor{lyred}{模}(长度). 
设
(
)
,
,
,
n
x
x x
x
=
\Lambda 
,
(
)
,
,
,
n
n
y
y y
y
=
\in 
\Lambda 
\mathbb{R},则两者之间的\textcolor{lyred}{距离}定义为 
(
)
(
)
(
)
(
)
,
n
n
x y
y
x
y
x
y
x
\rho 
=
-
+
-
+
+
-
\Lambda 
,也记为x
y
-
. 
利用距离概念可以定义邻域,由此得到开集,闭集,区域等概念. 
\vspace{4pt}
{\bfseries \textcolor{lyred}{二.多元函数的概念}}

\textcolor{lyred}{定义}.设D 为
\mathbb{R}上非空点集,若按照法则f ,使得
(
)
,
P x y
D
\forall 
\in 
,变量z 均有唯一 
确定的(实数)值与之对应,则称z 为x , y 的\textcolor{lyred}{二元函数},记为
(
)
,
z
f x y
=
,或者 
()
z
f P
=
,其中,x y 为\textcolor{lyred}{自变量}, z 为\textcolor{lyred}{因变量},或函数值. 
类似地,可定义三元及三元以上的\textcolor{lyred}{n 元函数}:
()
u
f P
=
,
n
P
D
\in 
\subset \mathbb{R}. 
\textcolor{lyred}{例}.二元函数
z
x
y
=
+
,它的图形为旋转抛物面. 
\vspace{4pt}
{\bfseries \textcolor{lyred}{三.多元函数的极限}}

\textcolor{lyred}{定义}.给定
()
f P ,设
0P 为
f
D 的聚点,若
\varepsilon 
\forall 
>
,
\delta 
\exists 
>
,当
(
)
0,
P
D
U P \delta 
\in 
\cap 
\circ 
时, 
()
f P
A
\varepsilon 
-
<
,则称A 为当
P
P
\to 
时,
()
f P 的\textcolor{lyred}{极限},记为
()
lim
P
P f P
A
\to 
=
,或者, 
(
)
(
)
(
)
,
,
lim
,
x y
x
y
f
x y
A
\to 
=
,也称为\textcolor{lyred}{二重极限}. 
\textcolor{lyred}{注}.
()
lim
P
P f P
\to 
存在的本质是当D 中动点P 以任何方式趋向
0P 时,
()
f P 均趋向于 
同一个常数;由于平面上
P
P
\to 
的方式有无限多种,故这个要求是相当高的. 
\textcolor{lyred}{例}.设
(
)
,
sin
xy
xy
f x y
x
y
x
y
=
+
+
,证明:
(
)
(
)
(
)
,
0,0
lim
,
x y
f x y
\to 
=
. 
证.
(
)
,
xy
f x y
x
y
x
y
-
\le 
\le 
+
+
,因此,
\varepsilon 
\forall 
>
,可取
\delta 
\varepsilon 
=
, 
当
(
)
(
)
x
y
\delta 
<
-
+
-
<
时,
(
)
,
f x y
\varepsilon 
-
<
,证毕. 
\textcolor{lyred}{例}.设
(
)
,
xy
f x y
x
y
=
+
,证明:
(
)
(
)
(
)
,
0,0
lim
,
x y
f x y
\to 
不存在. 
证.(
)
,x y 沿y
kx
=
趋向于(
)
0,0 时,
(
)
,
k
f x y
k
\to +
,依k 不同而不同,故 
(
)
(
)
,
0,0
lim
x y
xy
x
y
\to 
+
不存在,证毕. 
\textcolor{lyred}{注}.若沿不同的路径,当
P
P
\to 
时,
()
f P 趋向于不同的数,则
()
lim
P
P f P
\to 
不存在. 
\textcolor{lyred}{例}.设
(
)
,
xy
f x y
x
y
=
+
,证明:
(
)
(
)
(
)
,
0,0
lim
,
x y
f x y
\to 
不存在. 
证.当(
)
,x y 沿y
kx
=
趋向(
)
0,0 时,
(
)
,
k x
k x
f x y
x
k x
k x
=
=
\to 
+
+
; 
当(
)
,x y 沿
x
y
=
趋向(
)
0,0 时,
(
)
,
y
f x y
y
y
=
\to 
+
,证毕. 
\textcolor{lyred}{例}.
(
)
(
)
(
)
(
)
,
0,1
,
0,1
1 1
lim
lim
1 1
x y
x y
xy
xy
xy
\to 
\to 
+-=
=
++
. 
\textcolor{lyred}{例}.
(
)
(
)
(
)
(
)
(
)
(
)
(
)
(
)
,
1,0
,
1,0
,
1,0
sin
lim
lim
lim
x y
x y
x y
x y
x y
x
x
xy
x
xy
x
\to 
\to 
\to 
=
=
=
-
-
-
. 
\textcolor{lyred}{例}.求
(
)
(
)
(
)
,
0,0
lim
sin
x y
x
y
x y
x
y
\to 
+
+
. 
解.
(
)
(
)
,
2 sin
f x y
x y
x y
\le 
\le 
\le 
,故
(
)
(
)
(
)
,
0,0
lim
sin
x y
x
y
x y
x
y
\to 
+
=
+
. 
\textcolor{lyred}{例}.求
(
)
(
)
,
0,0
lim
x y
x
y
x
y
\to 
+
+
. 
解.
(
)
,
x
y
f x y
x
y
x
y
x
y
x
y
\le 
\cdots 
+
\cdots 
\le 
+
+
+
,故
(
)
(
)
,
0,0
lim
x y
x
y
x
y
\to 
+
=
+
. 
\textcolor{lyred}{例}.求
(
)
(
)
,
0,0
lim
x y
x y
x
y
\to 
+
. 
解.
(
)
(
)
,
0,0
cos
sin
lim
lim
lim
cos
sin
x y
x y
x
y
\rho 
\rho 
\rho 
\theta 
\theta 
\rho 
\theta 
\theta 
\rho 
+
+
\to 
\to 
\to 
=
=
=
+
. 
\vspace{4pt}
{\bfseries \textcolor{lyred}{四.多元函数的连续性}}

\textcolor{lyred}{定义}.给定
()
f P ,
f
P
D
\in 
,
0P 是
f
D 的聚点,若
()
(
)
lim
P
P f P
f P
\to 
=
,则称
0P 为
()
f P 的 
\textcolor{lyred}{连续点},
()
f P 在
0P 处\textcolor{lyred}{连续}. 
\textcolor{lyred}{定义}.若
()
f P 在开或者闭区域D (其中的点均是聚点)上处处连续,则称
()
f P 为 
D 上的\textcolor{lyred}{连续函数}. 
\textcolor{lyred}{定理}.多元连续函数的和差积商(分母不为零)以及复合仍连续. 
\textcolor{lyred}{定理}.多元初等函数在它的定义区域均连续. 
\textcolor{lyred}{定义}.由具有不同自变量的基本初等函数和常数函数经过有限多次四则运算以及 
复合所形成的,并且可以用一个式子来表示的函数称为\textcolor{lyred}{多元初等函数}. 
\textcolor{lyred}{例}.
(
)
(
)
2 sin
, ,
ln
x
yz
f x y z
e
x
z
+
=
+
+
+
是三元初等函数. 
\vspace{4pt}
{\bfseries \textcolor{lyred}{五.有界闭区域上连续函数的性质}}

\textcolor{lyred}{定理(最大值最小值定理)}.有界闭区域上的连续函数必有界,并且能在该区域上 
取到最大值和最小值. 
\textcolor{lyred}{定理(介值定理)}.有界闭区域上的连续函数能够取到介于最大值和最小值之间的 
任何值. 
\textcolor{lyred}{定义}.设
()
f P 为D 上的函数,若
\varepsilon 
\forall >
,
\delta 
\exists 
>
,当
,
P P
D
\in 
,且
PP
\delta 
<
时,均有 
(
)
(
)
f P
f P
\varepsilon 
-
<
,则称
()
f P 在D 上\textcolor{lyred}{一致连续}. 
\textcolor{lyred}{定理}.有界闭区域上的连续函数必在该区域上一致连续. 
\vspace{6pt}
{\large \bfseries \textcolor{lyred}{第9.2 节 偏导数}}

\vspace{4pt}
{\bfseries \textcolor{lyred}{一.偏导数的定义及其计算法}}

设
(
)
,
z
f x y
=
在(
)
,
x
y
的邻域内有定义,固定
y
y
=
,则()
(
)
,
g x
f x y
=
成为x 的 
一元函数,若
(
)
(
)
(
)
,
,
lim
x
f x
x y
f x y
g x
x
\Delta \to 
+\Delta 
-
'
=
\Delta 
存在,则称为
(
)
,
f x y 在(
)
,
x
y
处 
对x 的\textcolor{lyred}{偏导数},记为
(
)
,
xz
x
y
,
(
)
,
xf
x
y
,
x x
y y
z
x
=
=
\partial 
\partial 
,
x x
y y
f
x
=
=
\partial 
\partial 
,等价地, 
(
)
(
)
,
,
lim
x
x
x x
y y
f x y
f x
y
z
x
x
x
\to 
=
=
-
\partial 
=
\partial 
-
; 
类似地,
(
)
(
)
(
)
(
)
,
,
,
,
lim
lim
y
y
y
x x
y y
f x
y
y
f x
y
f x
y
f x
y
z
y
y
y
y
\Delta \to 
\to 
=
=
+\Delta 
-
-
\partial 
=
=
\partial 
\Delta 
-
.\textcolor{lyred}{ }
若
(
)
,
z
f
x y
=
在区域D 内处处存在偏导数,则得到它的\textcolor{lyred}{偏导函数 }
(
)
(
)
,
,
lim
x
f x
x y
f x y
z
x
y
\Delta \to 
+\Delta 
-
\partial =
\partial 
\Delta 
,
(
)
(
)
,
,
lim
y
f x y
y
f x y
z
y
y
\Delta \to 
+\Delta 
-
\partial =
\partial 
\Delta 
. 
\textcolor{lyred}{例}.设
z
x
xy
y
=
+
+
,则
z
x
y
x
\partial =
+
\partial 
,
z
x
y
y
\partial =
+
\partial 
. 
\textcolor{lyred}{例}.设
2 tan 2
z
x
y
=
,则
2 tan 2
z
x
y
x
\partial =
\partial 
,
sec 2
z
x
y
y
\partial =
\partial 
. 
\textcolor{lyred}{例}.设
(
)
sin
0,
y
z
x
x
x
=
>
\neq 
,则
sin
sin
y
z
y x
x
-
\partial =
\cdots 
\partial 
,
sin ln
cos
y
z
x
x
y
y
\partial =
\cdots 
\partial 
. 
\textcolor{lyred}{例}.理想气体状态方程PV
RT
=
,其中R 为常数,证明:
P
V
T
V
T
P
\partial 
\partial 
\partial 
\cdots 
\cdots 
=-
\partial 
\partial 
\partial 
. 
证.
RT
P
RT
P
V
V
V
\partial 
=
\Rightarrow 
=-
\partial 
,
RT
V
R
V
P
T
P
\partial 
=
\Rightarrow 
=
\partial 
,
PV
T
V
T
R
P
R
\partial 
=
\Rightarrow 
=
\partial 
,即得,证毕. 
\textcolor{lyred}{例}.设
(
)
(
)
,
1 arctan
y
y
f x y
x e
x
x
=
+
-
,求
(
)
1,0
xf
,
(
)
1,0
yf
. 
解.
(
)
,0
f x
x
=
,故
(
)
(
)
,0
1,0
x
x
f
x
x
f
=
\Rightarrow 
=
; 
(
)
1,
y
f
y
e
=
,故
(
)
(
)
1,
1,0
y
y
y
f
y
e
f
=
\Rightarrow 
=
. 
\textcolor{lyred}{例}.设
(
)
,
x
y
f x y
e
+
=
,求
(
)
0,0
xf
,
(
)
0,0
yf
. 
解.
(
)
,0
x
f x
e
=
,
(
)
0,0
lim
lim
lim
x
x
x
x
x
x
x
e
e
e
f
x
x
x
\to 
\to 
\to 
-
-
=
=
=
-
-
,不存在; 
(
)
0,
y
f
y
e
=
,
(
)
(
)
0,
0,0
y
y
y
f
y
ye
f
=
\Rightarrow 
=
. 
\textcolor{lyred}{例}.设
(
)
(
)
(
)
(
)
(
)
2 ,
,
0,0
,
0,
,
0,0
xy
x y
x
y
f x y
x y

\neq 

+
=

=

,求
(
)
0,0
xf
,
(
)
0,0
yf
. 
解.
(
)
(
)
(
)
,0
0,0
0,0
lim
lim
x
x
x
f x
f
f
x
x
\to 
\to 
-
-
=
=
=
-
,类似地, 
(
)
(
)
(
)
0,
0,0
0,0
lim
lim
y
y
x
f
y
f
f
y
y
\to 
\to 
-
-
=
=
=
-
. 
\textcolor{lyred}{注}.上述函数在(
)
0,0 处不连续,故二元函数的偏导数存在不能推出连续性. 
\vspace{4pt}
{\bfseries \textcolor{lyred}{二.高阶偏导数}}

\textcolor{lyred}{二阶偏导数}:(1)
z
z
x
x
x
\partial 
\partial 
\partial 


=


\partial 
\partial 
\partial 


,也记为
f
x
\partial 
\partial 
,
xx
f ,
xx
z ,也称\textcolor{lyred}{纯偏导}; 
(2)
2z
z
x y
y
x
\partial 
\partial 
\partial 


=


\partial \partial 
\partial 
\partial 


,也记为
2 f
x y
\partial 
\partial \partial ,
xy
f ,
xy
z ,也称\textcolor{lyred}{混合偏导}; 
(3)
2z
z
y x
x
y


\partial 
\partial 
\partial 
=


\partial \partial 
\partial 
\partial 


,混合偏导;(4)
z
z
y
y
y


\partial 
\partial 
\partial 
=


\partial 
\partial 
\partial 


,纯偏导. 
\textcolor{lyred}{例}.设
z
x y
xy
xy
=
-
-
+,则
xz
x y
y
y
=
-
-
,
yz
x y
xy
x
=
-
-, 
xx
z
xy
=
,
yy
z
x
xy
=
-
,
xy
z
x y
y
=
-
-,
yx
z
x y
y
=
-
-. 
\textcolor{lyred}{例}.设
(
)
(
)
(
)
(
)
(
)
2 ,
,
0,0
,
0,
,
0,0
x y
x y
x
y
f x y
x y

\neq 

+
=

=

,求
(
)
0,0
xy
f
,
(
)
0,0
yx
f
. 
解.
(
)
(
)
(
)
,0
0,0
0,0
lim
x
x
f x
f
f
x
\to 
-
=
=
-
,
(
)
(
)
(
)
0,
0,0
0,0
lim
y
y
f
y
f
f
y
\to 
-
=
=
-
, 
(
)
(
)
,
0,0
x y \neq 
时,
(
)
(
)
(
)
(
)
,
x
x y x
y
x y
x
x y
x y
f
x y
x
y
x
y
x
y
+
-
\cdots 
=
=
-
+
+
+
, 
(
)
(
)
(
)
(
)
,
y
x
x
y
x y
y
x
x y
f
x y
x
y
x
y
x
y
+
-
\cdots 
=
=
-
+
+
+
,故 
(
)
(
)
(
)
0,
0,0
0,0
lim
lim
x
x
xy
y
y
f
y
f
f
y
y
\to 
\to 
-
-
=
=
=
-
-
, 
(
)
(
)
(
)
,0
0,0
0,0
lim
lim
y
y
yx
x
x
f
x
f
x
f
x
x
\to 
\to 
-
-
=
=
=
-
-
. 
\textcolor{lyred}{定理}.设
(
)
,
z
f x y
=
的混合偏导数
2z
y x
\partial 
\partial \partial ,
2z
x y
\partial 
\partial \partial 均在区域D 内连续,则在该区域内 
z
z
y x
x y
\partial 
\partial 
=
\partial \partial 
\partial \partial . 
\textcolor{lyred}{注}.类似地,若
(
)
,
z
f x y
=
的三阶混合偏导
3z
y x x
\partial 
\partial \partial \partial ,
3z
x y x
\partial 
\partial \partial \partial ,
3z
x x y
\partial 
\partial \partial \partial 在区域D 内 
连续,则必在D 内相等,可以统一记为
z
x
y
\partial 
\partial 
\partial . 
\textcolor{lyred}{例}.验证
ln
z
x
y
=
+
满足\textcolor{lyred}{Laplace 方程}:
z
z
x
y
\partial 
\partial 
+
=
\partial 
\partial 
. 
证.
z
x
x
x
y
\partial =
\partial 
+
,
(
)
(
)
(
)
x
y
x
x
z
y
x
x
x
y
x
y
\cdots 
+
-\cdots 
\partial 
-
=
=
\partial 
+
+
; 
类似地,
(
)
z
x
y
y
x
y
\partial 
-
=
\partial 
+
,故
z
z
x
y
\partial 
\partial 
+
=
\partial 
\partial 
. 
\textcolor{lyred}{例}.验证
u
x
y
z
=
+
+
满足\textcolor{lyred}{Laplace 方程}:
u
u
u
x
y
z
\partial 
\partial 
\partial 
+
+
=
\partial 
\partial 
\partial 
. 
证.
u
r
=
,
u
r
x
x
x
r
x
r
r
r
\partial 
\partial 
=-
=-
\cdots 
=-
\partial 
\partial 
,
u
r
x
r
r
x
x
r
x
r
r
\partial 
-\cdots 
\cdots \partial \partial 
-
=-
=
\partial 
, 
类似地,
u
y
r
y
r
\partial 
-
=
\partial 
,
u
z
r
z
r
\partial 
-
=
\partial 
,故
u
u
u
x
y
z
\partial 
\partial 
\partial 
+
+
=
\partial 
\partial 
\partial 
,证毕. 
\textcolor{lyred}{补充练习 }
1.设
(
)
,ln
ln
x y
y
e
f x
y
x
x
x
x
-


-
=
-




,求
(
)
,
f x y 及
,
f
f
x
y
\partial 
\partial 
\partial 
\partial . 
解.设u
x
y
=
-
,
ln
v
x
=
,则
v
x
e
=
,
v
y
e
u
=
-
,故
(
)
,
u
v
u
f u v
e
v
-
=
,即 
(
)
,
x
y
x
f x y
e
y
-
=
,于是
x
y
f
x e
x
y
-
\partial 
+
=
\partial 
,
(
)
x
y
x
y
f
e
y
y
-
+
\partial 
=-
\partial 
. 
2.设
(
)
(
)
(
)
(
)
(
)
sin
,
,
0,0
,
0,
,
0,0
x
y
x y
x
y
f x y
x y

+
\neq 

+
=

=

,求
(
)
0,0
xf
,
(
)
0,0
yf
. 
解.
(
)
(
)
(
)
,0
0,0
0,0
lim
lim
sin
x
x
x
x
f x
f
f
x
x
x
\to 
\to 
-
=
=
-
不存在,
(
)
0,0
yf
类似. 
3.设
()
f u 二阶可导,且
(
)
sin
x
z
f e
y
=
满足
x
z
z
ze
x
y
\partial 
\partial 
+
=
\partial 
\partial 
,求
()
f u . 
解.令
sin
x
u
e
y
=
,则
()
sin
x
z
f
u e
y
x
\partial 
'
=
\partial 
,
()
cos
x
z
f
u e
y
y
\partial 
'
=
\partial 
, 
()
()
sin
sin
x
x
z
f
u e
y
f
u e
y
x
\partial 
''
'
=
+
\partial 
,
()
()
cos
sin
x
x
z
f
u e
y
f
u e
y
y
\partial 
''
'
=
-
\partial 
, 
故
()
()
x
x
f
u e
f u e
''
=
,即
()
z
f u
=
满足z
z
''=
,解得
()
u
u
f u
C e
C e-
=
+
. 
4.设
(
)
u
f
x
y
=
+
满足
u
u
u
u
x
y
x
y
x
x
\partial 
\partial 
\partial 
+
-
\cdots 
+
=
+
\partial 
\partial 
\partial 
,求
()
f x . 
解.令
r
x
y
=
+
,则u
df
r
x df
x
dr
x
r dr
\partial 
\partial 
=
=
\partial 
\partial 
,
u
r
x df
x d f
x
r
dr
r
dr
\partial 
-
=
+
\partial 
, 
类似地,
u
r
y df
y d f
y
r
dr
r
dr
\partial 
-
=
+
\partial 
,代入原方程,得
d f
f
r
dr
+
=
,解得 
()
cos
sin
f r
C
r
C
r
r
=
+
+
-,即
()
cos
sin
f x
C
x
C
x
x
=
+
+
-. 
5.设在区域D 内
(
)
,
xf
x y 与
(
)
,
yf
x y 均存在且有界,证明:
(
)
,
f x y 在D 内连续. 
证.
(
)
(
)
(
)
(
)
,
,
,
,
z
f x
x y
y
f x y
y
f x y
y
f x y
\Delta 
=
+\Delta 
+\Delta 
-
+\Delta 
+
+\Delta 
-
\le  
(
)
(
)
(
)
(
)
,
,
,
,
f x
x y
y
f x y
y
f x y
y
f x y
+\Delta 
+\Delta 
-
+\Delta 
+
+\Delta 
-
= 
(
)
(
)
,
,
x
y
f
y
y
x
f
x
y
M
x
N
y
\psi 
\eta 
+\Delta 
\Delta +
\Delta 
\le 
\Delta +
\Delta ,故(
)
(
)
,
0,0
x
y
\Delta 
\Delta 
\to 
时, 
(
)
(
)
,
,
f x
x y
y
f x y
+\Delta 
+\Delta 
\to 
,证毕. 
\vspace{6pt}
{\large \bfseries \textcolor{lyred}{第9.3 节 全微分}}

\vspace{4pt}
{\bfseries \textcolor{lyred}{一.全微分的定义}}

\textcolor{lyred}{定义}.若
(
)
,
z
f x y
=
在(
)
,x y 处的\textcolor{lyred}{全增量}
(
)
(
)
,
,
z
f x
x y
y
f x y
\Delta =
+\Delta 
+\Delta 
-
可表示为 
()
z
A
x
B
y
o \rho 
\Delta =
\cdots \Delta +
\cdots \Delta +
,其中A , B 不依赖于x
\Delta , y
\Delta ,而
(
)
(
)
x
y
\rho =
\Delta 
+\Delta 
,则 
称
(
)
,
z
f x y
=
在(
)
,x y 处\textcolor{lyred}{可微分}, A x
B y
\Delta +
\Delta 称为它在(
)
,x y 处的\textcolor{lyred}{全微分},记为dz . 
\textcolor{lyred}{例}.设z
xy
=
,则
(
)(
)
z
x
x
y
y
xy
x x
y y
x y
\Delta =
+\Delta 
+\Delta 
-
=\Delta +\Delta +\Delta \Delta ,而 
(
)
(
)
(
)
(
)
(
)
(
)
(
)
(
)
x y
x
y
x
y
x
y
x
y
\Delta \Delta 
\Delta 
+\Delta 
\le 
\cdots 
\le 
\Delta 
+\Delta 
\Delta 
+\Delta 
\Delta 
+\Delta 
,故
lim
x y
\rho 
\rho 
\to 
\Delta \Delta =
,可微,且 
dz
x x
y y
=\Delta +\Delta . 
\vspace{4pt}
{\bfseries \textcolor{lyred}{二.可微的条件}}

\textcolor{lyred}{定理}.设
(
)
,
f x y 在(
)
,x y 处可微,则它在该点处连续. 
\textcolor{lyred}{定理}.设
(
)
,
z
f x y
=
在(
)
,x y 处可微,则它在(
)
,x y 处的偏导数必存在,并且 
z
z
dz
x
y
x
y
\partial 
\partial 
=
\Delta +
\Delta 
\partial 
\partial 
,即
z
z
dz
dx
dy
x
y
\partial 
\partial 
=
+
\partial 
\partial 
. 
\textcolor{lyred}{注}.
(
)
,
z
f x y
=
在(
)
,x y 处可微
(
)
(
)
(
)
(
)
,
,
lim
x
y
z
f
x y
x
f
x y
y
x
y
\rho 
+
\to 


\Delta -
\Delta +
\Delta 


\Leftrightarrow 
=
\Delta 
+\Delta 
. 
\textcolor{lyred}{例}.设
(
)
(
)
(
)
(
)
(
)
2 ,
,
0,0
,
0,
,
0,0
xy
x y
x
y
f x y
x y

\neq 

+
=

=

,讨论它在(
)
0,0 处的可微性. 
解.
(
)
0,0
xf
=
,
(
)
0,0
yf
=
,
(
)
(
)
(
)
(
)
,
0,0
x y
f
f
x
y
f
x
y
\Delta \Delta 
\Delta =
\Delta 
\Delta 
-
=
\Delta 
+\Delta 
, 
(
)
(
)
(
)
(
)
0,0
0,0
lim
lim
x
y
f
f
x
f
y
x y
x
y
\rho 
\rho 
\rho 
+
+
\to 
\to 


\Delta -
\Delta +
\Delta 
\Delta \Delta 

=
\Delta 
+\Delta 
不存在,不可微. 
\textcolor{lyred}{例}.设
(
)
(
)
(
)
(
)
(
)
2 ,
,
0,0
,
0,
,
0,0
x y
x y
x
y
f x y
x y

\neq 

+
=

=

,讨论它在(
)
0,0 处的可微性. 
解.
(
)
0,0
xf
=
,
(
)
0,0
yf
=
,
(
)
(
)
(
)(
)
(
)
(
)
,
0,0
x
y
f
f
x
y
f
x
y
\Delta 
\Delta 
\Delta =
\Delta 
\Delta 
-
=
\Delta 
+\Delta 
, 
(
)
(
)
(
)(
)
(
)
(
)
(
)
(
)
0,0
0,0
lim
lim
x
y
f
f
x
f
y
x
y
x
y
x
y
\rho 
\rho 
\rho 
+
+
\to 
\to 


\Delta -
\Delta -
\Delta 
\Delta 
\Delta 

=
=


\Delta 
+\Delta 
\Delta 
+\Delta 


,可微. 
\textcolor{lyred}{定理}.设
(
)
,
z
f x y
=
在(
)
,x y 的邻域内可偏导,且偏导函数在(
)
,x y 处连续,则它在 
(
)
,x y 处可微. 
\vspace{4pt}
{\bfseries \textcolor{lyred}{三.全微分在近似计算中的应用}}

设
(
)
,
f x y 在(
)
,
x
y
处可微,且
(
)
,
xf
x
y
与
(
)
,
yf
x
y
不同时为零,则 
当(
)
(
)
x
x
y
y
-
+
-
<<时,有近似公式: 
(
)
(
)
(
)(
)
(
)(
)
,
,
,
,
x
y
f x y
f x
y
f
x
y
x
x
f
x
y
y
y
\approx 
+
-
+
-
. 
\textcolor{lyred}{例}.近似计算(
)
2.02
1.04
. 
解.设
(
)
,
y
f x y
x
=
,
(
)
(
)
(
)
(
)
1.04,2.02
1,2
1,2
0.04
1,2
0.02
x
y
f
f
f
f
\approx 
+
\cdots 
+
\cdots 
, 
(
)
(
)
1,2
1,2
y
xf
yx -
=
=
,
(
)
(
)
1,2
1,2
ln
y
yf
x
x
=
=
,故(
)
2.02
1.04
1.08
\approx 
. 
\textcolor{lyred}{课内练习 }
1.设
(
)
,
z
f x y
=
在点(
)
1,2 处的全微分
dz
dx
dy
=
+
,求
(
)
2 ,2
lim
x
f
h
h
h
\to 
+
-
. 
解.
(
)
1,2
xf
=
,
(
)
1,2
yf
=
,故 
(
)
(
)
(
)
(
)(
)
(
)
1,2
1,2
2 ,2
1,2
lim
lim
x
y
x
x
f
h
f
h
o
h
h
f
h
h
f
h
h
\to 
\to 
\cdots 
+
\cdots -
+
+
+
-
-
=
(
)
2 2
lim
x
h
h
h
\to 
\cdots 
+\cdots -
=
=-. 
2.设
(
)
(
)
sin
,
,
0,
x
y
x
y
x
y
f x y
x
y

+
+
\neq 

+
=

+
=

,证明:
(
)
,
f x y 在(
)
0,0 处可微. 
证.
(
)
0,0
xf
=
,
(
)
0,0
yf
=
,
(
)
(
)
(
)
0,0
0,0
0,0
lim
x
y
f
f
x
f
y
\rho 
\rho 
+
\to 


\Delta 
-
\Delta -
\Delta 

= 
(
)
(
)
(
)
(
)
(
)
(
)
lim
sin
x
y
x
y
x
y
\rho 
+
\to 
\Delta 
+\Delta 
=
\Delta 
+\Delta 
\Delta 
+\Delta 
,故可微,证毕. 
3.证明:
(
)
,
f
x y
x
y
=
+
在(
)
0,0 处不可微. 
证.
(
)
0,0
xf
=,
(
)
0,0
yf
=,
(
)
(
)
(
)
0,0
0,0
0,0
lim
x
y
f
f
x
f
y
\rho 
\rho 
+
\to 


\Delta 
-
\Delta -
\Delta 

= 
(
)
(
)
(
)
(
)
lim
x
y
x
y
x
y
\rho 
+
\to 
\Delta 
+\Delta 
-\Delta -\Delta 
\Delta 
+\Delta 
不存在,故不可微,证毕. 
4.设
(
)
(
)
,
,
f x y
x
y
x y
\varphi 
=
+
,其中(
)
,x y
\varphi 
在(
)
0,0 处连续,证明:
(
)
,
f x y 在 
(
)
0,0 处可微
(
)
0,0
\varphi 
\Leftrightarrow 
=
. 
证.必要性:在(
)
0,0 处可微
(
)
0,0
xf
\Rightarrow 
存在,而
(
)
(
)
,0
0,0
lim
x
f x
f
x
\to 
-
=
-
(
)
(
)
(
)
,0
lim
lim
,0
0,0
x
x
x
x
x
x
x
x
\varphi 
\varphi 
\varphi 
\to 
\to 
-
=
=\pm 
-
,故只能(
)
0,0
\varphi 
=
; 
充分性:
(
)
(
)
(
)
(
)
,
0,0
0,0
0,0
lim
x
y
f
x
y
f
f
x
f
y
\rho 
\rho 
+
\to 


\Delta 
\Delta 
-
-
\Delta -
\Delta 

= 
(
)
(
)
(
)
(
)
(
)
(
)
,
,
lim
lim
lim
,
0,0
x
y
x
y
f
x
y
x
y
\rho 
\rho 
\rho 
\varphi 
\varphi 
\varphi 
\rho 
\rho 
+
+
+
\to 
\to 
\to 
\Delta 
+\Delta 
\Delta 
\Delta 
\Delta 
\Delta 
=
=
\Delta 
\Delta 
=
=
,故 
在(
)
0,0 处可微,证毕. 
\vspace{6pt}
{\large \bfseries \textcolor{lyred}{第9.4 节 多元复合函数的求导法则}}

\vspace{4pt}
{\bfseries \textcolor{lyred}{一.多元链式法则}}

\textcolor{lyred}{定理}.设
(
)
,
z
f u v
=
有连续的偏导数,且
()
u
u x
=
,
()
v
v x
=
可导,则 
()
()
,
z
f u x
v x
=



可导,并且dz
f
du
f
dv
dx
u dx
v dx
\partial 
\partial 
=
\cdots 
+
\cdots 
\partial 
\partial 
. 
\textcolor{lyred}{注}.设
(
)
, ,
z
f u v w
=
,
()
u
u x
=
,
()
v
v x
=
,
()
w
w x
=
,则 
dz
f
du
f
dv
f
dw
dx
u dx
v dx
w dx
\partial 
\partial 
\partial 
=
\cdots 
+
\cdots 
+
\cdots 
\partial 
\partial 
\partial 
. 
\textcolor{lyred}{推论}.设
(
)
, ,
z
f u v x
=
有连续的偏导数,且
()
u
u x
=
,
()
v
v x
=
可导,则 
()
()
,
,
z
f u x
v x
x
=



可导,并且dz
f
du
f
dv
f
dx
u dx
v dx
x
\partial 
\partial 
\partial 
=
\cdots 
+
\cdots 
+
\partial 
\partial 
\partial . 
\textcolor{lyred}{例}.设
sin
z
uv
x
=
+
,其中
x
u
e
=
,
cos
v
x
=
,求dz
dx . 
解.
sin
cos
cos
sin
cos
x
x
x
dz
z du
z dv
z
ve
u
x
x
e
x
e
x
x
dx
u dx
v dx
x
\partial 
\partial 
\partial 
=
\cdots 
+
\cdots 
+
=
-
+
=
-
+
\partial 
\partial 
\partial 
. 
\textcolor{lyred}{定理}.设
(
)
,
z
f u v
=
有连续的偏导数,且
(
)
,
u
u x y
=
,
(
)
,
v
v x y
=
可偏导,则 
(
)
(
)
,
,
,
z
f u x y
v x y
=



可偏导,并且z
f
u
f
v
x
u
x
v
x
\partial 
\partial 
\partial 
\partial 
\partial 
=
\cdots 
+
\cdots 
\partial 
\partial 
\partial 
\partial 
\partial , z
f
u
f
v
y
u
y
v y
\partial 
\partial \partial 
\partial \partial 
=
+
\partial 
\partial 
\partial 
\partial \partial . 
\textcolor{lyred}{注}.设
(
)
, ,
z
f u v w
=
,
(
)
,
u
u x y
=
,
(
)
,
v
v x y
=
,
(
)
,
w
w x y
=
,则 
z
f
u
f
v
f
w
x
u
x
v
x
w
x
\partial 
\partial 
\partial 
\partial 
\partial 
\partial 
\partial 
=
\cdots 
+
\cdots 
+
\cdots 
\partial 
\partial 
\partial 
\partial 
\partial 
\partial 
\partial , z
f
u
f
v
f
w
y
u
y
v
y
w
y
\partial 
\partial 
\partial 
\partial 
\partial 
\partial 
\partial 
=
\cdots 
+
\cdots 
+
\cdots 
\partial 
\partial 
\partial 
\partial 
\partial 
\partial 
\partial 
. 
\textcolor{lyred}{推论}.设
(
)
, ,
u
f x y z
=
有连续的偏导数,且
(
)
,
z
z x y
=
可偏导,则 
(
)
, ,
,
u
f
x y z x y
=



可偏导,并且u
f
f
z
x
x
z
x
\partial 
\partial 
\partial 
\partial 
=
+
\cdots 
\partial 
\partial 
\partial 
\partial , u
f
f
z
y
y
z
y
\partial 
\partial 
\partial 
\partial 
=
+
\cdots 
\partial 
\partial 
\partial 
\partial . 
\textcolor{lyred}{例}.设
(
)
xy
z
x
y
=
+
,求z
x
\partial 
\partial , z
y
\partial 
\partial . 
解.设
v
z
u
=
,其中
u
x
y
=
+
,v
xy
=
,则z
z
u
z
v
x
u
x
v x
\partial 
\partial \partial 
\partial \partial 
=
+
=
\partial 
\partial 
\partial 
\partial \partial 
(
)
(
)
(
)
1 2
ln
ln 2
xy
xy
v
v
vu
u
u y
xy
x
y
x
y
x
y
y
-
-\cdots +
\cdots 
=
+
\cdots +
+
+
\cdots 
, 
(
)
(
)
(
)
ln 2
xy
xy
z
z
u
z
v
xy
x
y
x
y
x
y
x
y
u
y
v y
-
\partial 
\partial \partial 
\partial \partial 
=
+
=
+
\cdots +
+
+
\cdots 
\partial 
\partial 
\partial 
\partial \partial 
. 
\textcolor{lyred}{例}.设
(
)
, ,
x
y
z
u
f x y z
e
+
+
=
=
,其中
2 sin
z
x
y
=
,求u
x
\partial 
\partial , u
y
\partial 
\partial . 
解.
2 sin
x
y
z
x
y
z
u
f
f
z
xe
ze
x
y
x
x
z
x
+
+
+
+
\partial 
\partial 
\partial 
\partial 
=
+
\cdots 
=
+
\cdots 
\partial 
\partial 
\partial 
\partial 
, 
cos
x
y
z
x
y
z
u
f
f
z
ye
ze
x
y
y
y
z
y
+
+
+
+
\partial 
\partial 
\partial 
\partial 
=
+
\cdots 
=
+
\cdots 
\partial 
\partial 
\partial 
\partial 
. 
\textcolor{lyred}{例}.求
z
z
x
y
x
y
\partial 
\partial 
+
=
\partial 
\partial 
在极坐标
x
y
\rho =
+
,
arctan y
x
\theta =
下的形式. 
解.
(
)(
)
,
,
z
x y
\rho \theta 
-
-
,则
x
x
\rho 
\rho 
\partial 
=
\partial 
,
y
y
\rho 
\rho 
\partial 
=
\partial 
,
y
x
\theta 
\rho 
\partial 
-
=
\partial 
,
x
y
\theta 
\rho 
\partial 
=
\partial 
,故 
z
z
z
z
x
z
y
x
x
x
\rho 
\theta 
\rho 
\theta 
\rho \rho 
\theta 
\rho 
\partial 
\partial 
\partial 
\partial 
\partial 
\partial 
\partial 
-
=
\cdots 
+
\cdots 
=
\cdots 
+
\cdots 
\partial 
\partial 
\partial 
\partial 
\partial 
\partial 
\partial 
, 
z
z
z
z
y
z
x
y
y
y
\rho 
\theta 
\rho 
\theta 
\rho \rho 
\theta 
\rho 
\partial 
\partial 
\partial 
\partial 
\partial 
\partial 
\partial 
=
\cdots 
+
\cdots 
=
\cdots 
+
\cdots 
\partial 
\partial 
\partial 
\partial 
\partial 
\partial 
\partial 
,代入
z
z
x
y
x
y
\partial 
\partial 
+
=
\partial 
\partial 
,得
z
\rho \rho 
\partial 
=
\partial 
. 
\textcolor{lyred}{例}.求
z
z
x
y
\partial 
\partial 
+
=
\partial 
\partial 
在极坐标
x
y
\rho =
+
,
arctan y
x
\theta =
下的形式. 
解.
(
)(
)
,
,
z
x y
\rho \theta 
-
-
,
z
z
x
z
y
x
\rho \rho 
\theta 
\rho 
\partial 
\partial 
\partial 
-
=
\cdots 
+
\cdots 
\partial 
\partial 
\partial 
,
z
x
\partial 
=
\partial 
x
y
z
z
x
z
z
z
y
z
x
x
x
x
x
x
\rho 
\rho 
\rho 
\rho 
\rho 
\theta 
\rho 
\theta 
\rho 
\rho \theta 
\rho 
\rho 
\rho 
\theta \rho 
\theta 
\rho 
\theta 
\rho 
\partial 
\partial 
\cdots 
-\cdots 
\cdots 
\cdots 




\partial 
\partial 
\partial 
\partial 
\partial 
\partial 
\partial 
\partial 
\partial 
-
\partial 
\partial 
\partial 
+
\cdots 
+
\cdots 
+
+
\cdots 
+
\cdots 
=




\partial 
\partial 
\partial \partial 
\partial 
\partial 
\partial \partial 
\partial 
\partial 
\partial 
\partial 




z x
z
y
x
z
x
z
x
z
y
y
z
xy
\rho 
\rho 
\rho 
\rho \theta \rho 
\rho 
\rho 
\rho 
\theta \rho \rho 
\theta 
\rho 
\rho 
\theta 
\rho 




\partial 
\partial 
-
\partial 
-
\partial 
\partial 
-
-
\partial 
+
\cdots 
+
\cdots 
+
+
\cdots 
+
\cdots 
=




\partial 
\partial \partial 
\partial 
\partial \partial 
\partial 
\partial 




z x
z
xy
z y
z
x
z
xy
\rho 
\rho 
\rho 
\rho \theta \rho 
\theta 
\rho 
\rho 
\rho 
\theta 
\rho 
\partial 
\partial 
\partial 
\partial 
-
\partial 
-
+
+
\cdots 
+
\cdots 
\partial 
\partial \partial 
\partial 
\partial 
\partial 
, 
(
)(
)
,
,
z
x y
\rho \theta 
-
-
,
z
z
y
z
x
y
\rho \rho 
\theta 
\rho 
\partial 
\partial 
\partial 
=
\cdots 
+
\cdots 
\partial 
\partial 
\partial 
,
z
y
\partial 
=
\partial 
y
x
z
z
y
z
z
z
x
z
y
y
y
y
y
y
\rho 
\rho 
\rho 
\rho 
\rho 
\theta 
\rho 
\theta 
\rho 
\rho \theta 
\rho 
\rho 
\rho 
\theta \rho 
\theta 
\rho 
\theta 
\rho 
\partial 
\partial 
\cdots 
-\cdots 
-\cdots 
\cdots 




\partial 
\partial 
\partial 
\partial 
\partial 
\partial 
\partial 
\partial 
\partial 
\partial 
\partial 
\partial 
+
\cdots 
+
\cdots 
+
+
\cdots 
+
\cdots 
=




\partial 
\partial 
\partial \partial 
\partial 
\partial 
\partial \partial 
\partial 
\partial 
\partial 
\partial 




z y
z
x
y
z
y
z
y
z x
x
z
xy
\rho 
\rho 
\rho 
\rho \theta \rho 
\rho 
\rho 
\rho 
\theta \rho \rho 
\theta 
\rho 
\rho 
\theta 
\rho 




\partial 
\partial 
\partial 
-
\partial 
\partial 
\partial 
-
+
\cdots 
+
\cdots 
+
+
\cdots 
+
\cdots 
=




\partial 
\partial \partial 
\partial 
\partial \partial 
\partial 
\partial 




z y
z
xy
z x
z
y
z
xy
\rho 
\rho 
\rho 
\rho \theta \rho 
\theta 
\rho 
\rho 
\rho 
\theta 
\rho 
\partial 
\partial 
\partial 
\partial 
-
\partial 
-
+
+
+
\cdots 
+
\cdots 
\partial 
\partial \partial 
\partial 
\partial 
\partial 
, 
代入
z
z
x
y
\partial 
\partial 
+
=
\partial 
\partial 
,得
z
z
z
z
z
x
y
\rho 
\rho 
\theta 
\rho 
\rho 
\partial 
\partial 
\partial 
\partial 
\partial 
+
=
+
+
\partial 
\partial 
\partial 
\partial 
\partial 
. 
\textcolor{lyred}{例}.设
(
)
,
z
f ax
by cx
dy
=
+
+
,其中f 有连续的二阶偏导,求
2z
x y
\partial 
\partial \partial . 
解.令u
ax
by
=
+
,v
cx
dy
=
+
,则
u
v
u
v
z
u
v
f
f
af
cf
af
cf
x
x
x
\partial 
\partial 
\partial 
=
+
=
+
=
+
\partial 
\partial 
\partial 
, 
uu
uv
vu
vv
uu
uv
vu
vv
z
u
v
u
v
a
f
f
c
f
f
abf
adf
bcf
cdf
x y
y
y
y
y




\partial 
\partial 
\partial 
\partial 
\partial 
=
+
+
+
=
+
+
+
=




\partial \partial 
\partial 
\partial 
\partial 
\partial 




(
)
(
)
uu
uv
vv
abf
ad
bc f
cdf
abf
ad
bc f
cdf
+
+
+
=
+
+
+
. 
\textcolor{lyred}{例}.设
(
)
,
w
f x
y
z xyz
=
+
+
,其中f 有连续的二阶偏导,求
2w
x z
\partial 
\partial \partial . 
解.令u
x
y
z
=
+
+
,v
xyz
=
,则
w
f
u
f
v
f
yzf
x
u
x
v
x
\partial 
\partial 
\partial 
\partial 
\partial 
=
\cdots 
+
\cdots 
=
+
\partial 
\partial 
\partial 
\partial 
\partial 
, 
(
)
(
)
(
)
f
f
w
f
yzf
yf
yz
f
xyf
yf
yz f
xyf
x z
z
z
z
\partial 
\partial 
\partial 
\partial 
=
+
=
+
+
=
+
+
+
+
=
\partial \partial 
\partial 
\partial 
\partial 
(
)
f
y x
z f
xy zf
yf
+
+
+
+
. 
\vspace{4pt}
{\bfseries \textcolor{lyred}{二.全微分的形式不变性}}

设
(
)
,
z
f u v
=
,
(
)
,
u
u x y
=
,
(
)
,
v
v x y
=
,均有连续的偏导数,则 
z
z
z
u
z
v
z
u
z
v
z
z
dz
dx
dy
dx
dy
du
dv
x
y
u
x
v x
u
y
v y
u
v


\partial 
\partial 
\partial \partial 
\partial \partial 
\partial \partial 
\partial \partial 
\partial 
\partial 


=
+
=
+
+
+
=
+




\partial 
\partial 
\partial 
\partial 
\partial \partial 
\partial 
\partial 
\partial \partial 
\partial 
\partial 




,即 
无论u ,v 是中间变量还是自变量,
z
z
dz
du
dv
u
v
\partial 
\partial 
=
+
\partial 
\partial 
总是对的. 
\textcolor{lyred}{补充练习 }
1.设
(
)
,
f x y 有连续的偏导数,且
(
)
1,1
f
=,
(
)
1,1
xf
=
,
(
)
1,1
yf
=
,令 
()
(
)
,
,
F x
f
x f x x
=



,求
()1
F'
. 
解.
()
(
)
(
)
(
)
(
)
,
,
,
,
,
,
F
x
f
x f x x
f
x f x x
f
x x
f
x x
'
=
+
+









,故 
()
(
)
(
)
(
)
(
)
1,1
1,1
1,1
1,1
F
f
f
f
f
'
=
+
+
=




. 
2.设
()
()
,
z
f
x
y
y
x
\varphi 
\zeta 
=
-
+



,其中f 有连续的二阶偏导,\varphi 与\zeta 可导,求
2z
x y
\partial 
\partial \partial . 
解.
()
z
f
x
f
x
\varphi 
\partial 
'
'
'
=
\cdots 
+
\partial 
,
(
)
()
()
(
)
z
f
f
y
x
f
x y
\zeta 
\varphi 
\partial 
''
''
'
'
''
=
\cdots -
+
\cdots 
+
\cdots -
+




\partial \partial 
()
()
()
()
()
f
y
x f
y
x
f
y f
\zeta 
\varphi 
\zeta 
\varphi 
\zeta 
''
'
'
''
'
'
''
'
''
\cdots 
=-
+
-
+




. 
3.设
, x
y
z
f
xy
g
y
x




=
+








,其中f 有连续的二阶偏导, g 二阶可导,求
2z
x y
\partial 
\partial \partial . 
解.
z
y
y
f
y
f
g
yf
f
g
x
y
x
y
x
\partial 
-
'
'
=
\cdots 
+
\cdots 
+
\cdots 
=
+
-
\partial 
, 
z
x
x
y
f
y
f
x
f
f
f
x
f
g
g
x y
y
y
y
y
x
x
x




\partial 
-
-
'
''
=
+
\cdots +
\cdots 
-
+
\cdots +
\cdots 
-
-
\cdots 
=




\partial \partial 




x
y
f
f
xyf
f
g
g
y
y
x
x
'
''
-
+
-
-
-
. 
4.设
B
AC
-
>
,且
A \neq 
,证明:存在非奇异线性变换
x
y
x
y
\psi 
\lambda 
\eta 
\mu 
=
+


=
+

,使得方程 
u
u
u
A
B
C
x
x y
y
\partial 
\partial 
\partial 
+
+
=
\partial 
\partial \partial 
\partial 
可化为
u
\psi \eta 
\partial 
=
\partial \partial 
. 
证. u
u
u
u
u
x
x
x
\psi 
\eta 
\lambda 
\mu 
\psi 
\eta 
\psi 
\eta 
\partial 
\partial 
\partial 
\partial 
\partial 
\partial 
\partial 
=
+
=
+
\partial 
\partial 
\partial 
\partial 
\partial 
\partial 
\partial 
, u
u
u
u
u
y
y
y
\psi 
\eta 
\psi 
\eta 
\psi 
\eta 
\partial 
\partial 
\partial 
\partial 
\partial 
\partial 
\partial 
=
+
=
+
\partial 
\partial 
\partial 
\partial 
\partial 
\partial 
\partial 
, 
u
u
u
u
u
u
u
u
x
x
x
x
x
\psi 
\eta 
\psi 
\eta 
\lambda 
\mu 
\lambda 
\lambda \mu 
\mu 
\psi 
\psi \eta 
\eta \psi 
\eta 
\psi 
\psi \eta 
\eta 




\partial 
\partial 
\partial 
\partial 
\partial 
\partial 
\partial 
\partial 
\partial 
\partial 
\partial 
\partial 
=
+
+
+
=
+
+




\partial 
\partial 
\partial 
\partial \partial 
\partial 
\partial \partial \partial 
\partial 
\partial 
\partial 
\partial \partial 
\partial 




, 
(
)
u
u
u
u
u
u
u
u
x y
y
y
y
y
\psi 
\eta 
\psi 
\eta 
\lambda 
\mu 
\lambda 
\lambda 
\mu 
\mu 
\psi 
\psi \eta 
\eta \psi 
\eta 
\psi 
\psi \eta 
\eta 




\partial 
\partial 
\partial 
\partial 
\partial 
\partial 
\partial 
\partial 
\partial 
\partial 
\partial 
\partial 
=
+
+
+
=
+
+
+




\partial \partial 
\partial 
\partial 
\partial \partial 
\partial 
\partial \partial 
\partial 
\partial 
\partial 
\partial 
\partial \partial 
\partial 




, 
u
u
u
u
u
u
u
u
y
y
y
y
y
\psi 
\eta 
\psi 
\eta 
\psi 
\psi \eta 
\eta \psi 
\eta 
\psi 
\psi \eta 
\eta 
\partial 
\partial 
\partial 
\partial 
\partial 
\partial 
\partial 
\partial 
\partial 
\partial 
\partial 
\partial 
=
+
+
+
=
+
+
\partial 
\partial 
\partial 
\partial \partial 
\partial 
\partial \partial 
\partial 
\partial 
\partial 
\partial 
\partial \partial 
\partial 
,于是 
(
)
(
)
(
)
u
u
u
A
B
C
A
B
C
A
B
C
\lambda 
\lambda 
\lambda \mu 
\lambda 
\mu 
\mu 
\mu 
\psi 
\psi \eta 
\eta 
\partial 
\partial 
\partial 
+
+
+
+
+
+
+
+
+
=




\partial 
\partial \partial 
\partial 
, 
故取\lambda , \mu 为
Ax
Bx
C
+
+
=
的两个不同的实根即可,证毕. 
\vspace{6pt}
{\large \bfseries \textcolor{lyred}{第9.5 节 隐函数的求导公式}}

\vspace{4pt}
{\bfseries \textcolor{lyred}{一.一个方程的情形}}

\textcolor{lyred}{定理}.设
(
)
,
F x y 在(
)
,
P x
y
的邻域内有连续偏导数,
(
)
,
F x y
=
,
(
)
,
yF x y
\neq 
, 
则方程
(
)
,
F x y =
在P 的某个邻域内唯一确定连续可导的函数
()
y
y x
=
,它满足 
()
,
F x y x
\equiv 




,
(
)
y
y x
=
,并且
x
y
F
dy
dx
F
=-
. 
\textcolor{lyred}{定理}.设
(
)
, ,
F x y z 在
(
)
,
,
P x
y z
的某个邻域内有连续的偏导数,
(
)
,
,
F x y z
=, 
(
)
,
,
zF
x
y z
\neq 
,则
(
)
, ,
F x y z =
在P 的某个邻域内唯一确定连续可偏导的函数 
(
)
,
z
z x y
=
,它满足
(
)
, ,
,
F x y z x y
\equiv 




,
(
)
,
z
z x
y
=
,并且
x
z
F
z
x
F
\partial =-
\partial 
,
y
z
F
z
y
F
\partial =-
\partial 
. 
\textcolor{lyred}{例}.设
(
)
,
z
z x y
=
是
x
y
z
z
+
+
-
=
确定的隐函数,求
z
x
\partial 
\partial 
. 
解.
x
z
F
z
x
x
F
z
\partial =-
=
\partial 
-
,
(
)
(
)
(
)
(
)
(
)
x
z
x
z
z
x
z
x
z
z
-
-\cdots -
-
+
\partial 
=
=
\partial 
-
-
. 
\textcolor{lyred}{例}.设
(
)
,
z
z x y
=
是
y z
x
ze +
=
确定的隐函数,求
2z
x y
\partial 
\partial \partial . 
解.两边对x 求偏导
(
)
(
)
y z
y z
y z
z
z
z
z
e
ze
x
x
x
z e
x
z
+
+
+
\partial 
\partial 
\partial 
\Rightarrow =
+
\Rightarrow 
=
=
\partial 
\partial 
\partial 
+
+
; 
两边对y 求偏导
(
)
y z
y z
y z
y z
z
z
z
ze
z
e
ze
y
y
y
z e
z
+
+
+
+


\partial 
\partial 
\partial 
-
\Rightarrow 
=
+
+
\Rightarrow 
=
=-


\partial 
\partial 
\partial 
+
+


, 
(
)
(
)
yz
z
z
z
x y
x
y
z
x
y
z
x
z
x
z
\partial 
\partial 
\partial 
-




=
\cdots 
=
\cdots 
=
=-




\partial \partial 
\partial 
+
\partial 
+




+
+
. 
\vspace{4pt}
{\bfseries \textcolor{lyred}{二.方程组的情形}}

\textcolor{lyred}{定理}.设
(
)
(
)
, , ,
, , ,
F x y u v
G x y u v

在
(
)
,
,
,
P x
y u v
的某个邻域内具有连续的偏导数,且 
(
)
(
)
,
,
,
,
,
,
F x
y u v
G x
y u v
=

=

,\textcolor{lyred}{Jacobi 行列式}(
)
(
)
,
,
u
v
u
v
F
F
F G
G
G
u v
\partial 
=
\partial 
在P 处不为0,则方程组 
(
)
(
)
, , ,
, , ,
F x y u v
G x y u v
=

=

在P 的邻域内唯一确定一组有连续偏导数的函数
(
)
(
)
,
,
u
u x y
v
v x y
=
=

, 
它满足
(
)
(
)
(
)
(
)
, ,
,
,
,
, ,
,
,
,
F x y u x y
v x y
G x y u x y
v x y

\equiv 






\equiv 






,
(
)
(
)
,
,
u
u x
y
v
v x
y
=

=

,且 
(
)
(
)
(
)
(
)
,
,
,
,
F G
x v
F G
u v
u
x
\partial 
\partial 
\partial 
\partial 
\partial 
=-
\partial 
,
(
)
(
)
(
)
(
)
,
,
,
,
F G
u x
F G
u v
v
x
\partial 
\partial 
\partial 
\partial 
\partial =-
\partial 
,
(
)
(
)
(
)
(
)
,
,
,
,
F G
y v
F G
u v
u
y
\partial 
\partial 
\partial 
\partial 
\partial 
=-
\partial 
,
(
)
(
)
(
)
(
)
,
,
,
,
F G
u y
F G
u v
v
y
\partial 
\partial 
\partial 
\partial 
\partial =-
\partial 
. 
\textcolor{lyred}{例}.设一组函数
(
)
(
)
,
,
u
u x y
v
v x y
=
=

,满足
xu
yv
yu
xv
-
=


+
=

,求u
x
\partial 
\partial , u
y
\partial 
\partial , v
x
\partial 
\partial , v
y
\partial 
\partial . 
解.方程组两边对x 求偏导,
u
v
x
y
u
x
x
u
v
y
x
v
x
x
\partial 
\partial 

-
=-
\partial 
\partial 
\partial 
\partial 

+
=-

\partial 
\partial 

,解得
u
xu
yv
x
x
y
v
yu
xv
x
x
y
\partial 
-
-

=
\partial 
+
\partial 
-

=
\partial 
+

; 
方程组两边对y 求偏导,
u
v
x
y
v
y
y
u
v
y
x
u
y
y
\partial 
\partial 

-
=
\partial 
\partial 
\partial 
\partial 

+
=-

\partial 
\partial 

,解得
u
xv
yu
y
x
y
v
xu
yv
y
x
y
\partial 
-

=
\partial 
+
\partial 
-
-

=
\partial 
+

. 
\textcolor{lyred}{例}.设
(
)
(
)
,
,
u
u x y
v
v x y
=
=

是
cos
sin
v
x
u
u
v
y
u
u
=


=

的反函数组,求u
x
\partial 
\partial , u
y
\partial 
\partial , v
x
\partial 
\partial , v
y
\partial 
\partial . 
解.
cos
sin
cos
sin
sin
sin
cos
sin
cos
cos
x
x
u
v
v
v
v
v
v
u
v
v
u
x
u
u
u
u
u
u
x
u
x
u
v
v
v
v
v
v
u
v
v
u
x
u
u
u
u
u
u
x
u
x

\partial 

\partial 
\partial 




=
-
\cdots 
+
-
\cdots 
=






\partial 
\partial 
\partial 





\Rightarrow 


\partial 
\partial 
\partial 






=
+
\cdots 
-
+
\cdots 
=






\partial 
\partial 
\partial 






,得 
cos
cos
sin
u
v
x
u
v
v
v
v
x
u
u
u
\partial 

=
\partial 
\partial 

=
-
\partial 

,类似地,
sin
sin
cos
u
v
y
u
v
v
v
v
y
u
u
u
\partial 

=
\partial 
\partial 

=
+
\partial 

. 
\textcolor{lyred}{定理}.设
(
)
(
)
, ,
, ,
F x y z
G x y z

在(
)
,
,
P x
y z
的邻域内有连续的偏导数,
(
)
(
)
,
,
,
,
F x
y
z
G x
y
z
=

=

, 
(
)
(
)
,
,
P
F G
y z
\partial 
\neq 
\partial 
,则方程组
(
)
(
)
, ,
, ,
F x y z
G x y z
=

=

在P 的邻域内唯一确定一组连续可导的 
函数
()
()
y
y x
z
z x
=
=

,它满足
()
()
()
()
,
,
,
,
F x y x
z x
G x y x
z x

\equiv 






\equiv 






,
(
)
(
)
y
y x
z
z x
=

=

,并且 
(
)
(
)
(
)
(
)
,
,
,
,
F G
F G
dy
x z
y z
dx
\partial 
\partial 
=-\partial 
\partial 
,
(
)
(
)
(
)
(
)
,
,
,
,
F G
F G
dz
y x
y z
dx
\partial 
\partial 
=-\partial 
\partial 
. 
\textcolor{lyred}{例}.设
(
)
,
y
f x t
=
,而
(
)
,
t
t x y
=
是由
(
)
, ,
F x y t =
确定,求dy
dx . 
解一.
(
)
,
,
x
t
x
t
t
t
x
t
y
f
f
dy
t
t dy
dy
y
f
x t x y
f
f
dx
x
y dx
dx
f
\partial 
\partial 
\partial 
\partial 
+


\partial 
\partial 
\equiv 
\Rightarrow 
=
+
+
\Rightarrow 
=






\partial 
\partial 
-


,而 
(
)
, ,
x
t
F
t
F x y t
x
F
\partial 
=
\Rightarrow 
=-
\partial 
,
y
t
F
t
y
F
\partial =-
\partial 
,代入,得
x
t
t
x
t
t
y
f F
f F
dy
dx
F
f F
-
=
+
; 
解二.
(
)
(
)
,
, ,
dy
f
f dt
dy
f dt
f
y
f x t
dx
x
t dx
dx
t dx
x
F
F dy
F dt
F dy
F dt
F
F x y t
x
y dx
t dx
y dx
t dx
x
\partial 
\partial 
\partial 
\partial 


=
+
-
=


=

\partial 
\partial 
\partial 
\partial 



\Rightarrow 
\Rightarrow 



\partial 
\partial 
\partial 
\partial 
\partial 
\partial 
=



+
+
=
+
=-
\partial 
\partial 
\partial 
\partial 
\partial 
\partial 




,解得 
x
t
t
x
t
t
y
f F
f F
dy
dx
F
f F
-
=
+
. 
\textcolor{lyred}{补充练习 }
1.设
(
)
,
f x y 有连续的偏导数,
y
x
=
是
(
)
,
f x y 的一条等高线,已知
(
)
1,1
yf
=-, 
求
(
)
1,1
xf
. 
解. (
)
(
)
(
)
,
,
,
x
y
f x x
C
f
x x
f
x x
x
=
\Rightarrow 
+
\cdots 
=
,代入
x =,
y =,得 
(
)
(
)
1,1
1,1 3
x
y
f
f
+
\cdots =
,故
(
)
1,1
xf
=
. 
2.设
(
)
,
z
f x
y
z xyz
=
+
+
,
(
)
,
f x y 有连续的偏导数,求z
x
\partial 
\partial , x
y
\partial 
\partial , y
z
\partial 
\partial . 
解.
x
z
F
f
yzf
z
x
F
f
xyf
+
\partial =-
=
\partial 
-
-
,
y
x
F
f
xzf
x
y
F
f
yzf
+
\partial =-
=-
\partial 
+
,
z
y
F
f
xyf
y
z
F
f
xzf
-
-
\partial =-
=
\partial 
+
. 
3.设
(
)
ln
z
f x
y
xy
=
-
+
,其中
()
y
y x
=
是
xy
e
x
y
-
+
=
确定的隐函数,已知 
()
()
f
f '
=
=,求
x
dz
dx
=
. 
解.
(
)
(
)
f
x
y
dz
dy
dy
y
x
dx
f x
y
dx
dx
'
-


=
-
+
+


-


,而
x =
时
y =-,代入得 
()
()
1 0
x
x
x
f
dz
dy
dy
dx
f
dx
dx
=
=
=
'


=
-
-+
=-




,由
xy
e
x
y
-
+
=
,两边求导,得 
xy
xy
xy
x
dy
dy
dy
ye
dy
e
y
x dx
dx
dx
xe
dx
=
-


+
-+
=
\Rightarrow 
=
\Rightarrow 
=


+


,故
x
dz
dx
=
=-. 
4.设
u
xy z
=
,其中
(
)
,
z
z x y
=
满足
x
y
z
xyz
+
+
=
,求
(
)
1,1,1
u
y
\partial 
\partial 
. 
解.
(
)
(
)
1,1,1
1,1,1
u
z
u
z
xyz
xy z
y
y
y
y
\partial 
\partial 
\partial 
\partial 
=
+
\Rightarrow 
=
+
\partial 
\partial 
\partial 
\partial 
,由
x
y
z
xyz
+
+
=
, 
两边求偏导,
(
)
1,1,1
z
z
z
y
xz
z
y
z
xz
xy
y
y
y
xy
z
y
\partial 
\partial 
\partial 
-
\partial 
+
=
+
\Rightarrow 
=
\Rightarrow 
=-
\partial 
\partial 
\partial 
-
\partial 
,故 
(
)
1,1,1
u
y
\partial 
=
-
=-
\partial 
. 
5.设
(
)
,
z
z x y
=
由方程
(
)
,
F x
y z
-
=
确定,其中
(
)
,
F u v 具有连续的二阶偏导数, 
求
2z
x y
\partial 
\partial \partial . 
解.令
(
)
(
)
, ,
,
f x y z
F x
y z
=
-
,则
x
z
f
F
z
x
f
F
\partial =-
=-
\partial 
,
y
z
f
F
z
y
f
F
\partial =-
=
\partial 
, 
(
)
(
)
z
z
F
F
F
F
F
F
y
y
F
z
z
x y
y
F
F




\partial 
\partial 
-
+
-
-
+






\partial 
\partial 
\partial 
\partial 




=
-
=-


\partial \partial 
\partial 

(
)
(
)
F
F
F
F
F
F
F
F
F
F
F




-
+
\cdots 
-
-
+
\cdots 








=-
[
]
[
]
F F
F F F
F
F F
F F
F F
F F F
F F
F
F
-
+
-
-
+
-
+
=-
=
. 
\vspace{6pt}
{\large \bfseries \textcolor{lyred}{第9.6 节 多元函数微分学的几何应用}}

\vspace{4pt}
{\bfseries \textcolor{lyred}{一.一元向量值函数及其导数}}

{\bfseries \textcolor{lyred}{1.基本概念}}

\textcolor{lyred}{定义}.设I \subset \mathbb{R},则
:
f
I \to 
\rho 
\mathbb{R}称为\textcolor{lyred}{(一元)向量值函数},记为
()
r
f t
=
\rho 
\rho 
,其中 
()
()
()
()
(
)
,
,
f t
x t
y t
z t
=
\rho 
, ()
x t , ()
y t , ()
z t 为其\textcolor{lyred}{分量}. 
\textcolor{lyred}{定义}.
()
()
()
()
(
)
lim
lim
,lim
,lim
t
t
t
t
t
t
t
t
f t
x t
y t
z t
\to 
\to 
\to 
\to 
=
\rho 
. 
\textcolor{lyred}{定义}.若
()
()
lim
t
t f t
f t
\to 
=
\rho 
\rho 
,则称
()
f t
\rho 
在0t 处\textcolor{lyred}{连续}. 
\textcolor{lyred}{定义}.
(
)
()
()
()
()
(
)
lim
lim
,
,
t
t
f t
t
f t
r
x t
y t
z t
t
t
\Delta \to 
\Delta \to 
+\Delta 
-
\Delta 
'
'
'
=
=
\Delta 
\Delta 
\rho 
\rho 
\rho 
,称为
()
r
f t
=
\rho 
\rho 
的\textcolor{lyred}{导数},或 
\textcolor{lyred}{导向量},记为dr
dt
\rho 
, df
dt
\rho 
,
()
r t
'\rho 
,
()
f
t
'
\rho 
. 
{\bfseries \textcolor{lyred}{2.几何意义}}

设
()
()
()
()
(
)
,
,
r
r t
x t
y t
z t
=
=
\rho 
\rho 
,若r
OM
=
\xi \xi \xi \xi \rho 
\rho 
,则M 的轨迹是一条曲线\Gamma ,称为()
r t
\rho 
的 
\textcolor{lyred}{终端曲线},或\textcolor{lyred}{图形},即
()
()
()
:
x
x t
y
y t
z
z t
=


\Gamma 
=


=

,而
()
r
f t
=
\rho 
\rho 
称为\Gamma 的\textcolor{lyred}{向量方程}. 
设()
r t
OM
=
\xi \xi \xi \xi \rho 
\rho 
, (
)
r t
t
OM'
+\Delta =
\xi \xi \xi \xi \rho 
\rho 
,则
r
MM
t
t
'
\Delta 
=
\Delta 
\Delta 
\xi \xi \xi \xi \xi \rho 
\rho 
为MM '的方向向量,而
lim
t
dr
r
dt
t
\Delta \to 
\Delta 
=
\Delta 
\rho 
\rho 
为 
\Gamma 在M 处切线的方向向量(假设不为零向量),称为\Gamma 的\textcolor{lyred}{切向量},记为T
\rho 
. 
{\bfseries \textcolor{lyred}{3.物理意义}}

设
()
r
f t
=
\rho 
\rho 
为动点M 的位置向量,则它的速度
dr
v
dt
=
\rho 
\rho 
,速率
dr
v
dt
=
\rho 
; 
加速度
dv
d r
a
dt
dt
=
=
\rho 
\rho 
\rho 
,加速度大小
d r
a
dt
=
\rho 
. 
{\bfseries \textcolor{lyred}{4.求导法则}}

(1)
()
()
()
()
f t
g t
f
t
g t
\alpha 
\beta 
\alpha 
\beta 
'


'
'
+
=
+


\rho 
\rho 
\rho 
\rho 
;(2)
()()
()()
()
()
f t g t
f
t g t
f t g t
'
'
'
=
+




\rho 
\rho 
\rho 
; 
(3)
()
()
()
()
()
()
f t
g t
f
t
g t
f t
g t
'


'
'
\cdots 
=
\cdots 
+
\cdots 


\rho 
\rho 
\rho 
\rho 
\rho 
\rho 
; 
(4)
()
()
()
()
()
()
f t
g t
f
t
g t
f t
g t
'


'
'
\times 
=
\times 
+
\times 


\rho 
\rho 
\rho 
\rho 
\rho 
\rho 
; 
(5)
()
\{
\}
()
()
f
g t
g t f
g t
'
'
'
=








\rho 
\rho 
. 
\vspace{4pt}
{\bfseries \textcolor{lyred}{二.空间曲线的切线与法平面}}

\textcolor{lyred}{情况1}.\textcolor{lyred}{参数方程} 
设空间曲线
()
()
()
:
x
x t
y
y t
z
z t
=


\Gamma 
=


=

,
(
)
,
,
M x
y z
对应参数
t
t
=
,则该点处的\textcolor{lyred}{切线}方程为: 
()
()
()
x
x
y
y
z
z
x t
y t
z t
-
-
-
=
=
'
'
'
,过M 且与切线垂直的平面称为\Gamma 在该点处的\textcolor{lyred}{法平面}. 
\textcolor{lyred}{例}.求
:
x
t
y
t
z
t
=


\Gamma 
=

=

在(
)
1,1,1 处的切线与法平面方程. 
解.
(
)
(
)
1,2 ,3
1,2,3
t
T
t
t
=
=
=
\rho 
,故切线方程为
x
y
z
-
-
-
=
=
; 
法平面方程为
(
)
(
)
(
)
x
y
z
\cdots 
-
+\cdots 
-
+\cdots 
-
=
. 
\textcolor{lyred}{情况2}.\textcolor{lyred}{一般方程} 
\textcolor{lyred}{定理}.设
(
)
(
)
, ,
:
, ,
F x y z
G x y z
=

\Gamma 
=

,若在M 处(
)(
)
,
,
,
,
x
y
z
x
y
z
F F F
G G G
\times 
\neq 
\rho 
,则\Gamma 在M 处的 
切向量
(
)(
)
,
,
,
,
x
y
z
x
y
z
T
F F F
G G G
=
\times 
\rho 
. 
\textcolor{lyred}{例}.求
:
xyz
x
y
=-

\Gamma -
+
=

垂直于2
x
y
z
=
=
的法平面方程. 
解.设切点(
)
, ,
x y z ,则
(
)(
)(
)
,
,
1, 2 ,0
,
, 2
T
yz xz xy
y
xy
xy
y z
xz
=
\times 
-
=
-
-
\rho 
,由 
(
)
// 2,1,1
xy
xy
y z
xz
T
y
-
-
\Rightarrow 
=
=
\Rightarrow 
=
\rho 
,由
xyz
x
x
y
z
=-
=-


\Rightarrow 


-
+
=
=


,故 
法平面方程为(
)(
)(
)
x
y
z
+
+
-
+
-
=
,即2
x
y
z
+
+
=
. 
\vspace{4pt}
{\bfseries \textcolor{lyred}{三.曲面的切平面与法线}}

在
(
)
:
, ,
F x y z
\Sigma 
=
上任取过
(
)
, ,
M x y z 的光滑曲线
()
: r
r t
\Gamma 
=
\rho 
\rho 
,则 
()
()
()
,
,
F x t
y t
z t
=




,两边对t 求导,得到(
)
,
,
,
,
x
y
z
dx dy dz
F F F
dt
dt
dt


\cdots 
=




,其中 
,
,
dx dy dz
T
dt
dt
dt


=



\xi \rho 
是\Gamma 的切向量,而
(
)
,
,
x
y
z
n
F F F
=
\rho 
与\Gamma 无关,于是\Sigma 上所有过M 的 
曲线在M 处的切线均垂直于n\rho ,它们构成一个平面,称为\Sigma 的\textcolor{lyred}{切平面},而n\rho 为它的 
法向量,称为曲面\Sigma 的\textcolor{lyred}{法向量};过M ,垂直于切平面的直线称为该曲面的\textcolor{lyred}{法线}. 
\textcolor{lyred}{例}.求
z
x
y
=
+
-在(
)
2,1,4 处的切平面及法线方程. 
解.
x
y
z
+
--
=
,
(
)(
)
(
)
2,1,4
2 ,2 , 1
4,2, 1
n
x
y
=
-
=
-
\rho 
,故切平面方程为 
(
)
(
)(
)
x
y
z
-
+
-
-
-
=
;法线方程为
x
y
z
-
-
-
=
=-
. 
\textcolor{lyred}{例}.求
x
y
z
+
+
=
上平行于
x
y
z
+
+
=
的切平面方程. 
解.设切点为(
)
, ,
x y z ,则
(
)(
)
2 ,4 ,6
// 1,4,6
x
y
z
n
x
y
z
=
\Rightarrow 
=
=
\rho 
,得2x
y
z
=
=
, 
代入
x
y
z
+
+
=
,得
x =\pm ,故切点为(
)
1, 2, 2
\pm \pm 
\pm 
,切平面方程为
(
)
(
)
(
)
x
y
z
\pm 
+\cdots 
\pm 
+\cdots 
\pm 
=
,即
x
y
z
+
+
=\pm 
. 
\textcolor{lyred}{例}.求过
:
x
y
z
L
x
y
z
+
-
-
=

+
-
=

,与
x
y
z
+
-
=
相切的平面方程. 
解.设切点
(
)
,
,
P x
y z
,则
x
y
z
+
-
=
,由
(
)
,
,
n
x
y
z
=
-
\rho 
,得切平面方程 
(
)
(
)
(
)
x
x
x
y
y
y
z
z
z
-
+
-
-
-
=
,即
x x
y y
z z
+
-
-
=
; 
设切平面为
(
)
x
y
z
x
y
z
\lambda 
+
-
-+
+
-
=
,则 
x
y
x
y
z
z
\lambda 
\lambda 
\lambda 
\lambda 
\lambda 
\lambda 
=
+

+
+
--
-

=
=
=
\Rightarrow 
=+

-
-

=+

,代入
x
y
z
+
-
=
,解得
\lambda =-, 7
-, 
即
(
)
1,0,0
P
或(
)
1, 6, 6
--
-
,故切平面方程为
x -=
,或
x
y
z
+
-
+=
. 
\textcolor{lyred}{例}.求过
:
x
y
z
L
-
-
-
=
=
-
,与
x
y
z
+
+
=
相切的平面方程. 
解.设切点
(
)
,
,
P x
y z
,则
x
y
z
+
+
=
,由
(
)
,2
,3
n
x
y
z
=
\rho 
,得切平面方程 
(
)
(
)
(
)
x
x
x
y
y
y
z
z
z
-
+
-
+
-
=
,即
x x
y y
z z
+
+
-
=
; 
代入
6,3, 2






,得
x
y
z
+
+
-
=
; 
最后,(
)(
)
2,1, 1
,2
,3
x
y
z
x
y
z
-
\cdots 
=
\Rightarrow 
+
-
=
,得切点为 
(
)
3,0,2 ,或(
)
1,2,2 ,故切平面为
x
y
z
+
+
-
=
,或
x
z
+
-
=
. 
\vspace{4pt}
{\bfseries \textcolor{lyred}{四.全微分的几何意义}}

曲面
(
)
:
,
z
f x y
\Sigma 
=
在点(
)
,
,
x
y
z
处的切平面方程为 
(
)(
)
(
)(
)(
)(
)
,
,
x
y
f
x
y
x
x
f
x
y
y
y
z
z
-
+
-
+-
-
=
,即 
(
)(
)
(
)(
)
,
,
x
y
z
z
f
x
y
x
x
f
x
y
y
y
-
=
-
+
-
,故
(
)
,
z
f x y
=
的全微分是曲面 
(
)
:
,
z
f x y
\Sigma 
=
的切平面上点的纵坐标增量. 
\textcolor{lyred}{补充练习 }
1.设
:
x
t
y
t
z
t
=


\Gamma 
=-

=

上与
ax
by
cz
d
+
+
+
=
平行的切线恰有二条,求, ,
a b c 应满足的 
条件. 
解.
(
)
(
)
(
)(
)
1, 2 ,3
, ,
1, 2 ,3
, ,
T
t
t
a b c
t
t
a b c
a
bt
ct
=
-
\bot
\Rightarrow 
-
\cdots 
=
\Rightarrow 
-
+
=
\rho 
,此方程 
恰有两个不同实根,故
b
ac
\Delta =
-
>
,且
c \neq 
. 
2.求过
x
y
z
x
y
z

+
+
=

+
+
=

在
1,
,
M 







处的切线,且垂直于
x
y
z
+
+
=
的 
平面方程. 
解.
(
)
(
)
2 ,2 ,2
2,
2, 2
M
n
x
y
z
=
=
\rho 
,
(
)
(
)
2 ,4 ,8
2,2 2,4 2
M
n
x
y
z
=
=
\rho 
,得 
(
)
2 2, 3 2, 2
T
n
n
=
\times 
=
-
\rho 
\rho 
\rho 
,故
(
)(
)
(
)
2, 3 2,
2,1,1
2,0, 2
n =
-
\times 
=-
-
\rho 
, 
于是所求平面方程为
(
)
x
y
z




\cdots 
-
+\cdots 
-
-\cdots 
-
=












,即
x
z
-
=
. 
3.设
z
x
y
=
+
在(
)
1,2,5 处的切平面通过
:
x
y
b
L
ax
y
z
+
=


+
-
=

,求,a b. 
解.
(
)(
)
(
)
1,2,5
2 ,2 , 1
2,4, 1
n
x
y
=
-
=
-
\rho 
,设
(
)
:
x
y
b
ax
y
z
\lambda 
\Pi 
+
-
+
+
-
-
=
,则 
a
b
\lambda 
\lambda 
\lambda 
\lambda 
\lambda 
+
+
-
--
=
=
=
\Rightarrow 
=
-
-
,
a =-,
b =
. 
\vspace{6pt}
{\large \bfseries \textcolor{lyred}{第9.7 节 方向导数与梯度}}

\vspace{4pt}
{\bfseries \textcolor{lyred}{一.方向导数}}

设
(
)
cos
:
cos
x
x
t
l
t
y
y
t
\alpha 
\beta 
=
+

\ge 

=
+

是以
0P 为始点,以\alpha , \beta 为方向角的射线,若
(
)
,
f x y 在 
某个
(
)
U P 内有定义,则
(
)
,
f x y 在
0P 处沿方向l 的\textcolor{lyred}{方向导数}为 
()
(
)
(
)
(
)
(
)
cos
, 
cos
,
lim
lim
P
P
l
t
P
f P
f P
f x
t
y
t
f x
y
f
l
PP
t
\alpha 
\beta 
+
\to 
\to 
-
+
+
-
\partial 
=
=
\partial 
沿着
. 
\textcolor{lyred}{例}.求
(
)
,
f x y
x
y
=
+
在
(
)
0,0
O
处沿
(
)
cos
,cos
e
\alpha 
\beta 
=
\rho 
的方向导数. 
解.
(
)
(
)
(
)
0,0
cos
, cos
0,0
lim
t
f t
t
f
f
l
t
\alpha 
\beta 
+
\to 
-
\partial 
=
=
\partial 
. 
\textcolor{lyred}{定理}.设
(
)
,
f x y 在(
)
,
P x y 处可微,则它在该点沿任意方向l 的方向导数均存在, 
并且
grad
l
f
f e
l
\partial 
=
\cdots 
\partial 
\rho ,其中grad
,
f
f
f
x
y


\partial 
\partial 
=

\partial 
\partial 


,
(
)
cos ,cos
le
\alpha 
\beta 
=
\rho 
. 
\textcolor{lyred}{例}.求
2 y
z
xe
=
在(
)
1,0
P
处沿(
)
1,0
P
到
(
)
2, 1
Q
-
方向的方向导数. 
解.
(
)(
)(
)
(
)
1,0
grad
1,0
,2
1,2
y
y
z
e
xe
=
=
,
(
)
1, 1
l =
-
\rho 
,
,
le


=
-




\rho 
,故 
(
)
(
)
1,0
1,2
,
z
l
\partial 


=
\cdots 
-
=-


\partial 


. 
\textcolor{lyred}{例}.求
(
)
, ,
f x y z
xy
yz
zx
=
+
+
在(
)
1,1,2 处沿l
\rho 
的方向导数,其中l
\rho 
的方向角分别为 
60\circ , 45\circ ,60\circ . 
解.
(
)
(
)(
)
(
)
1,1,2
grad
1,1,2
,
,
3,3,2
f
y
z x
z x
y
=
+
+
+
=
, 
(
)
2 1
cos60 ,cos45 ,cos60
,
,
le


=
=





\circ 
\circ 
\circ 
\rho 
,
(
)
1,1,2
3 2
grad
l
f
f e
l
\partial 
+
=
\cdots 
=
\partial 
\rho 
. 
\textcolor{lyred}{例}.求
x
y
u
z
+
=
在(
)
1,1,1
P
处沿曲面
x
y
z
+
+
=
在该点处指向内侧的 
法线方向的方向导数. 
解.
(
)(
)
(
)
1,1,1
4 ,6 ,2
4,6,2
n
x
y
z
=-
=-
\rho 
,故
,
,
ne
-
-
-


=



\rho 
,而 
grad
,
,
,
,
P
P
u
u
u
u
x
y
z


\partial 
\partial 
\partial 


=
=
-




\partial 
\partial 
\partial 




,故
grad
n
P
P
u
u
e
n
\partial 
=
\cdots 
=-
\partial 
\rho 
. 
\vspace{4pt}
{\bfseries \textcolor{lyred}{二.梯度}}

设
(
)
,
f x y 在(
)
,x y 的某个邻域内有连续偏导数,记
(
)
grad
,
,
f
f
f x y
x
y


\partial 
\partial 
=

\partial 
\partial 


,称为 
(
)
,
f x y 在(
)
,x y 处的\textcolor{lyred}{梯度},也记为
f
\nabla 
\rho 
,其中\nabla 
\rho 
称为\textcolor{lyred}{向量微分算子}. 
设
(
)
, ,
f x y z 在(
)
, ,
x y z 的某个邻域内有连续偏导数,则它在(
)
, ,
x y z 处的\textcolor{lyred}{梯度}为 
(
)
grad
, ,
,
,
f
f
f
f x y z
x
y
z


\partial 
\partial 
\partial 
=

\partial 
\partial 
\partial 


,也记为
f
\nabla 
\rho 
. 
\textcolor{lyred}{几何意义}:曲线
(
)
,
f x y
c
=
称为
(
)
,
f x y 的\textcolor{lyred}{等值线},grad f 是该曲线的法向量,并且 
从数值较低的等值线指向数值较高的等值线,沿该方向函数的变化率最大. 
曲面
(
)
, ,
f x y z
c
=
称为
(
)
, ,
f x y z 的\textcolor{lyred}{等值面},grad f 是该曲面的法向量,并且从数值 
较低的等值面指向数值较高的等值面,沿该方向函数的变化率最大. 
\textcolor{lyred}{定理}.设
(
)
,
f x y 在(
)
,x y 处可微,则 
(1)当
(
)
grad
,
l
f x y
=
\rho 
时,
(
)
grad
,
f
f x y
l
\partial 
=
\partial 
最大; 
(2)当
(
)
grad
,
l
f x y
=-
\rho 
时,
(
)
grad
,
f
f x y
l
\partial 
=-
\partial 
最小; 
(3)当
(
)
grad
,
l
f x y
\bot
\rho 
时,
f
l
\partial 
=
\partial 
. 
\textcolor{lyred}{例}.求
(
)
(
)
,
f x y
x
y
=
+
在(
)
1,1 处方向导数的最大值与最小值,以及使得方向 
导数为零的方向. 
解.
(
)
(
)(
)
(
)
1,1
grad
1,1
,
1,1
f
x y
=
=
,故f
l
\partial 
\partial 的最大值为
2 ,最小值为
-
; 
沿
(
)
1, 1
l =\pm 
-
\rho 
方向的方向导数为零. 
\textcolor{lyred}{例}.求
(
)
, ,
f x y z
x
xy
z
=
-
-
在(
)
1,1,0 处方向导数的最大值以及取到该最大值的 
方向. 
解.
(
)(
)(
)
(
)
1,1,0
grad
1,1,0
, 2
, 1
2, 2, 1
f
x
y
xy
=
-
-
-
=
-
-
,故当
(
)
2, 2, 1
l =
-
-
\rho 
时, 
(
)
2, 2, 1
f
l
\partial 
=
-
-
=
\partial 
最大. 
\vspace{4pt}
{\bfseries \textcolor{lyred}{三.数量场与向量场}}

\textcolor{lyred}{定义}.区域G 上的函数
(
)
f M 称为G 上的\textcolor{lyred}{数量场},如温度场;区域G 上的向量值 
函数
(
)
F M
\rho 
称为G 上的\textcolor{lyred}{向量场},如力场. 
对于数量场
(
)
f M ,向量场
(
)
grad
F
f M
=
\rho 
称为它的\textcolor{lyred}{梯度场},反之
(
)
f M 称为这个 
向量场的\textcolor{lyred}{势函数},具有势函数的向量场称为\textcolor{lyred}{保守场},\textcolor{lyred}{守恒场},\textcolor{lyred}{位势场},例如引力场, 
真空中的静电场. 
\textcolor{lyred}{例}.证明:中心力场
()
(
)
r
r
F
f r e
f
x
y
z
e
=
=
+
+
\rho 
\rho 
\rho 一定是保守场. 
证.设
()
()
G r
f r
'
=
,则
()
()
()
( , , )
grad
grad
r
x y z
G
G r
r
G r
f r e
r
'
'
=
=
=
\rho ,证毕. 
\textcolor{lyred}{补充练习 }
1.求
arctan
y
x
u
y
z


=




在(
)
1,1,1 处的最大方向导数. 
解. (
)
(
)
,1,1
1,1,1
x
u x
x
u
\pi 
\pi 
=
\Rightarrow 
=
, (
)
(
)
1, ,1
1,1,1
y
u
y
y
u
=
\Rightarrow 
=
, 
(
)
(
)
1,1,
1,1,1
z
u
z
z
u
\pi 
\pi 
-
=
\Rightarrow 
=-
,故
max
,2,
u
l
\pi 
\pi 
\pi 
\partial 


=
-
=
+


\partial 


. 
2.求由
z
xz
y
+
-
=
确定的
(
)
,
z
z x y
=
在(
)
0,1 处的最大方向导数. 
解.
x =
,
y =时
z =,
(
)
(
)
0,1
1 1
grad
0,1
,
,
3 3
z
z
z
x
z
x
-




=
=-




+
+




,故 
(
)
max
grad
0,1
z
z
l
\partial =
=
\partial 
. 
3.设
(
)
,
f x y
x
y
=
+
在
(
)
,
M x
y
处沿
(
)
1,2
l =
\rho 
方向有最大的增长率,并且 
(
)
,
f x
y
=
,求
(
)
,
M x
y
. 
解.
(
)
grad
1,2
M
f
y
=
与(
)
1,2 同向
y
\Rightarrow 
=,又
(
)
,
f x
y
x
=
\Rightarrow 
=
. 
4.设
(
)
(
)
,
1 ln
ax
f x y
ye
b x
y
=
+
+
在(
)
0,1 处沿
(
)
1,1
l =
\rho 
方向取得最大增长率2 2 , 
求a ,b . 
解.
(
)
(
)
grad
0,1
,1
f
a
b
=
+
与(
)
1,1 同向,故
a
b
a
b
+
=
\Rightarrow 
=
+,且
a >
,又 
(
)
,
2 2
a b
a
a
a
+
=
\Rightarrow 
+
=
\Rightarrow 
=
,
b =. 
5.设xOy 面上的一条曲线通过(
)
0,1 ,并且平行于
(
)
,
f x y
x
xy
=
-
的梯度场,求 
该曲线的方程. 
解.
(
)(
)
grad
,
, 6
x
y
f
f
f
x
y
xy
=
=
-
-
,设曲线方程为
()
y
y x
=
,于是 
dy
xy
dx
x
y
-
=
-
,此为齐次方程,解得
3x y
y
C
-
=
,
C =-. 
\vspace{6pt}
{\large \bfseries \textcolor{lyred}{第9.8 节 多元函数的极值及其求法}}

\vspace{4pt}
{\bfseries \textcolor{lyred}{一.多元函数的极值}}

设
(
)
,
f x y 在(
)
,
x
y
某个邻域内有定义,若对该邻域内不同于(
)
,
x
y
的点上均有
(
)
(
)
,
,
f x y
f x
y
<
,则称
(
)
,
f x y 在(
)
,
x
y
处取得\textcolor{lyred}{极大值};\textcolor{lyred}{极小值}类似. 
极大极小值统称为\textcolor{lyred}{极值},使函数取得极值的点称为\textcolor{lyred}{极值点}. 
\textcolor{lyred}{例}.设
(
)
,
f x y 连续,
(
)
(
)
(
)
(
)
,
0,0
,
lim
x y
f x y
xy
x
y
\to 
-
=
+
,证明:
(
)
0,0
f
不是极值. 
证.
(
)
(
)
(
)
,
f x y
xy
x
y
o
x
y


=
+
+
+
+



,故
(
)
(
)
,
f x x
x
o x
=
+
,而 
(
)
(
)
,
f x
x
x
o x
-
=-
+
,而
(
)
0,0
f
=
,即得,证毕. 
\textcolor{lyred}{定理(必要条件)}.设
(
)
,
f x y 在点(
)
,
x
y
处偏导数存在,且有极值,则 
(
)
,
xf
x
y
=
,
(
)
,
yf
x
y
=
. 
\textcolor{lyred}{定理(充分条件)}.设
(
)
,
f x y 在(
)
,
x
y
的某个邻域内有连续的二阶偏导数,且 
(
)
,
x
y
是其驻点,令
(
)
,
xx
A
f
x
y
=
,
(
)
,
xy
B
f
x
y
=
,
(
)
,
yy
C
f
x
y
=
,则 
(1)
AC
B
-
>
时(
)
,
x
y
为极值点,且
A <
时为极大值,
A >
时为极小值; 
(2)
AC
B
-
<
时(
)
,
x
y
不是极值点. 
\textcolor{lyred}{例}.求
(
)
,
f x y
x
y
x
y
x
=
-
+
+
-
的极值. 
解.
(
)
(
)
,
,
x
y
f
x y
x
x
f
x y
y
y

=
+
-
=

=-
+
=

,解得驻点(
)
1,0 ,(
)
1,2 ,(
)
3,0
-
,(
)
3,2
-
, 
xx
A
f
x
=
=
+
,
xy
B
f
=
=
,
yy
C
f
y
=
=-
+
,列表如下 
(
)
,x y        (
)
1,0        (
)
1,2       (
)
3,0
-
       (
)
3,2
-
A           12          12         12
-
-
B           0           0           0           0 
C           6          6
-           6          6
- 
AC
B
-
        72        72
-
-
(
)
,
f x y       极小值     非极值        非极值       极大值 
极小值
(
)
1,0
f
=-,极大值
(
)
3,2
f -
=
. 
\vspace{4pt}
{\bfseries \textcolor{lyred}{二.多元函数的最大值最小值}}

\textcolor{lyred}{方法}:(1)确定
(
)
,
f x y 在D 内最大最小值的存在性; 
(2)列出
(
)
,
f x y 在D 内的驻点以及偏导数不存在的点,计算这些点处的函数值; 
(3)若D 不是开区域,求出
(
)
,
f x y 在D
D
\partial 
\cap 
上的最大最小值; 
(4)比较上述值的大小,最大者为最大值,最小者为最小值. 
\textcolor{lyred}{例}.求体积为a 的长方体的最小表面积. 
解.设长宽为x , y ,高为a
xy ,故表面积
a
a
S
xy
x
y
xy
xy


=
+
\cdots 
+
\cdots 
=




(
)
0,
a
a
xy
x
y
y
x


+
+
>
>




,由
x
a
S
y
x


=
-
=




,
y
a
S
x
y


=
-
=




,得 
x
y
a
=
=
,唯一驻点,由实际情况, (
)
3,
S
a
a
a
=
为最小值. 
\textcolor{lyred}{例}.求
(
)
(
)
,
f x y
x y
x
y
=
-
-
在由
x
y
+
=
, x 轴和y 轴所围成的闭区域D 上的 
最大最小值. 
解. D 为有界闭区域,故
(
)
,
f x y 在D 上存在最大最小值; 
由
(
)
(
)
(
)
(
)
(
)
(
)
,
,
x
y
f
x y
xy
x
y
x y
xy
x
y
x
y
f
x y
x
x
y
x y
x
x
y
x
y

=
-
-
-
=
-
-
=
+
=


\Rightarrow 


=
-
-
-
=
-
-
=
+
=



,得 
唯一驻点(
)
2,1 ,而
(
)
2,1
f
=
; 在边界
x =
和
y =
上,
(
)
,
f x y =
; 
在
x
y
+
=
上,
(
)
,
f x y
x
x
=-
+
,0
x
\le 
\le 
,最小值为
(
)
4,2
f
=-
,最大值为 
(
)
(
)
0,6
6,0
f
f
=
=
;故最大值为4 ,最小值为64
-
. 
\vspace{4pt}
{\bfseries \textcolor{lyred}{三.条件极值 拉格朗日乘数法}}

\textcolor{lyred}{条件极值}:求目标函数
(
)
,
f x y 在约束条件(
)
,
x y
\varphi 
=
下的极值. 
\textcolor{lyred}{拉格朗日乘数法}.设(
)
,
x
y
为上述条件极值问题的解,则存在\lambda ,使得 
(
)
(
)
(
)
(
)
,
,
,
,
x
x
y
y
f
x
y
x
y
f
x
y
x
y
\lambda \varphi 
\lambda \varphi 
+
=

+
=

,其中\lambda 称为\textcolor{lyred}{拉格朗日乘子}. 
\textcolor{lyred}{注}.令
(
)
(
)
,
,
L
f x y
x y
\lambda \varphi 
=
+
,称为\textcolor{lyred}{拉格朗日函数},则上述条件为
x
y
L
L
L\lambda 
=
=
=
. 
\textcolor{lyred}{注}.其它情况,例如求
(
)
, , ,
f x y z t 在约束
(
)
(
)
, , ,
, , ,
x y z t
x y z t
\varphi 
\zeta 
=

=

下的极值,则令 
(
)
(
)
(
)
(
)
, , , ; ,
, , ,
, , ,
, , ,
L x y z t
f x y z t
x y z t
x y z t
\lambda \mu 
\lambda \varphi 
\mu \zeta 
=
+
\cdots 
+
\cdots 
,由 
x
y
z
t
L
L
L
L
=
=
=
=
及
\varphi 
\zeta 
=
=
,解出(
)
, ,
x y z 即可. 
\textcolor{lyred}{例}.求
(
)
,
f x y
x
y
=
+
在(
)
(
)
x
y
-
+
-
=
上的最大最小值. 
解.令(
)
(
)
(
)
, ;
L x y
x
y
x
y
\lambda 
\lambda 

=
+
+
-
+
-
-



,则 
(
)
(
)
(
)
(
)
2 2
x
y
L
x
x
L
y
y
x
y
\lambda 
\lambda 

=
+
\cdots 
-
=


=
+
\cdots 
-
=



-
+
-
=

,解得
x
y
=
\pm 

=
\pm 

,而 
(
)
3,
6 6
f
\pm 
\pm 
=
\pm 
,故最大值15
6 6
+
,最小值15
6 6
-
. 
\textcolor{lyred}{例}.求周长为12的直角三角形的最大面积. 
解.设直角边长为x , y ,则
S
xy
=
,令
(
)
L
xy
x
y
x
y
\lambda 
=
+
+
+
+
-
,则 
x
y
x
L
y
x
y
y
L
x
x
y
x
y
x
y
\lambda 
\lambda 






=
+
+
=



+








=
+
+
=



+




+
+
+
=



,由(1)(2),得
x
x
y
y
x
y
x
y
x
y
+
+
=
\Rightarrow 
=
+
+
,故 
6 2
x
y
=
=
-
,于是
max
72 2
A
=
-
. 
\textcolor{lyred}{例}.求表面积为
a 的长方体的最大体积. 
解.设长宽高为, ,
x y z ,则
xyz
V =
,令
(
)
L
xyz
xy
yz
zx
a
\lambda 
=
+
+
+
-
,则 
(
)
(
)
(
)
x
y
z
L
yz
y
z
L
xz
x
z
x
y
z
a
L
xy
y
x
xy
yz
xz
a
\lambda 
\lambda 
\lambda 

=
+
+
=

=
+
+
=

\Rightarrow 
=
=
=

=
+
+
=


+
+
=

,故
max
V
a
=
. 
\textcolor{lyred}{例}.求曲面
z
x
y
=
+
到平面
:
x
y
z
\Pi 
+
+
+=
的最短距离. 
解.
x
y
z
d
+
+
+
=
,令
(
)
(
)
L
x
y
z
z
x
y
\lambda 
=
+
+
+
+
-
-
,则 
(
)
(
)
(
)
x
y
z
L
x
y
z
x
L
x
y
z
y
x
y
L
x
y
z
z
x
y
\lambda 
\lambda 
\lambda 

=
+
+
+
-
=

=
+
+
+
-
=

\Rightarrow 
=
=-

=
+
+
+
+
=

=
+

,
z =
,故
min
d
=
. 
\textcolor{lyred}{例}.求原点到曲线
z
x
y
x
y
z
=
+
+
+
=

上点的最长与最短距离. 
解.令
(
)
(
)
L
x
y
z
z
x
y
x
y
z
\lambda 
\mu 
=
+
+
+
-
-
+
+
+
-
,由 
x
x
x
y
y
y
z
z
x
y
z
x
y
z
\lambda 
\mu 
\lambda 
\mu 
\lambda 
\mu 

-+
-
+
=

=



-
+
=


-+


\Rightarrow 
=
+
+
=




=
+

=
-
+
+
=




或
x
y
z

--
=



--

=


=
+


,故
max
min
5 3
5 3
d
d

=
+


=
-

. 
\textcolor{lyred}{补充练习 }
1.将正数a 分解成三个正数x , y , z 之和,使得
m
n
p
u
x y z
=
最大,其中m ,n , p 为 
已知正数. 
解.令
(
)
m
n
p
L
x y z
x
y
z
a
\lambda 
=
+
+
+
-
,则
m
n
p
x
m
n
p
y
m
n
p
z
L
mx
y z
L
nx y
z
L
px y z
x
y
z
a
\lambda 
\lambda 
\lambda 
-
-
-

=
+
=

=
+
=

=
+
=

+
+
=

,解得x
y
z
m
n
p
=
=
, 
故
m
x
a
m
n
p
=
+
+
,
n
y
a
m
n
p
=
+
+
,
p
z
a
m
n
p
=
+
+
. 
2.欲做无盖长方体容器,底部造价每平米3元,侧面每平米1元,设预算为36元,求 
尺寸使得容积最大. 
解.设长宽高为x , y , z ,令
(
)
1 2
L
xyz
xy
xz
yz
\lambda 
=
+
+\cdots 
+
-



,则 
(
)
(
)
(
)
(
)
1 2
x
y
z
L
yz
y
z
y
y
z
x
y
L
xz
x
z
x
x
z
z
x
z
L
xy
x
y
z
y
y
x
y
xy
xz
yz
\lambda 
\lambda 
\lambda 
=
+
+
=

+

=
=



=
+
+
=
+



\Rightarrow 
\Rightarrow 



+
=
+
+
=
=



=

+


+\cdots 
+
=


,代入 
(
)
1 2
xy
xz
yz
+\cdots 
+
=
,得
x =
,
y =
,
z =
. 
3.设长方体的三个面在坐标面上,一个顶点在平面
(
)
, ,
x
y
z
a b c
a
b
c
+
+
=
>
上,求 
该长方体的最大体积. 
解.设长方体在平面上的那个顶点为
(
)
, ,
A x y z ,则V
xyz
=
,令 
x
y
z
L
xyz
a
b
c
\lambda 

=
+
+
+
-




,则
x
y
z
L
yz
a
a
x
L
xz
b
b
y
L
xy
c
c
z
x
y
z
a
b
c
\lambda 
\lambda 
\lambda 

=
+
\cdots 
=

=




=
+
\cdots 
=


\Rightarrow 
=




=
+
\cdots 
=

=



+
+
=

,故
max
abc
V
=
. 
4.求过
2,1, 3






的平面
(
)
, ,
x
y
z
a b c
a
b
c
+
+
=
>
,使它与三个坐标面所围四面体的 
体积最小. 
解.设该平面方程为
x
y
z
a
b
c
+
+
=,其中2
a
b
c
+
+
=,体积
V
abc
=
, 
令
L
abc
a
b
c
\lambda 

=
+
+
+
-




,由
a
b
c
L
L
L
L\lambda 
=
=
=
=
,解得
a =
,
b =
,
c =, 
此时体积最小,故所求方程为
x
y
z
+
+
=. 
5.在曲面
x
y
z
+
+
=上求一点,使得函数
u
x
y
z
=
+
+
在该点沿方向 
(
)
1, 1,0
l =
-
\rho 
的方向导数最大. 
解.
(
)
grad
2 ,2 ,2
u
x
y
z
=
,
,
,0
le


=
-




\rho 
,
grad
l
u
u e
x
y
l
\partial 
=
\cdots 
=
-
\partial 
\rho 
, 
令
(
)
L
x
y
x
y
z
\lambda 
=
-
+
+
+
-
,则
x
y
z
L
x
L
y
L
z
x
y
z
\lambda 
\lambda 
\lambda 
=+
=


=-+
=

=
=


+
+
=

,解得 
x =\pm 
,
y =\mu 
,
z =
,故
,
,0
u
l 

-




\partial 
=
\partial 
最大,
1 1
, ,0
2 2
u
l 

-




\partial 
=-
\partial 
最小. 